% ------------------------------------------------------------
% English working translation layer.
% Source file: papers_30_59/37_Sur_la_variation_par_torsion_des_constantes_locales_d_equations_fonctionnelles_d/37_Constantes_Locales.tex
% Modern title: On the variation under twisting of local constants in functional equations of L-functions
% ------------------------------------------------------------

% Sur la variation, par torsion, des constantes locales d'equations fonctionnelles de fonctions L
% P. Deligne et G. Henniart
% Published in Invent. math. 64, 89-118 (1981)
% Typeset from the original PDF with faithful transcription

\documentclass[12pt,a4paper]{article}
\usepackage{fontspec}
\usepackage[english]{babel}
\usepackage{amsmath, amssymb, amsthm, amscd}
\usepackage{mathtools}
\usepackage{enumitem}
\usepackage{graphicx}
\usepackage{xcolor}
\usepackage[margin=2.5cm]{geometry}
\usepackage{hyperref}
\usepackage{tikz-cd}

% Theorem environments
\theoremstyle{plain}
\newtheorem{theorem}{Th\'eor\`eme}[section]
\newtheorem{lemma}[theorem]{Lemma}
\newtheorem{proposition}[theorem]{Proposition}
\newtheorem{corollary}[theorem]{Corollary}

\theoremstyle{definition}
\newtheorem{definition}[theorem]{D\'efinition}
\newtheorem{remark}[theorem]{Remark}

\theoremstyle{remark}
\newtheorem{example}[theorem]{Example}

% Math operators
\DeclareMathOperator{\Ind}{Ind}
\DeclareMathOperator{\Res}{Res}
\DeclareMathOperator{\Hom}{Hom}
% duplicate dim removed
% duplicate det removed
\DeclareMathOperator{\Tr}{Tr}
\DeclareMathOperator{\sw}{Sw}

% Shortcuts
\newcommand{\eps}{\varepsilon}
\newcommand{\cO}{\mathcal{O}}
\newcommand{\cX}{\mathcal{X}}
\newcommand{\bQ}{\mathbb{Q}}
\newcommand{\bZ}{\mathbb{Z}}
\newcommand{\bC}{\mathbb{C}}
\newcommand{\bR}{\mathbb{R}}

\let\dim\relax\let\det\relax

% --- Pass-2 compatibility shim. ---
% Purpose: repair transcription-level LaTeX omissions without changing mathematical text.
\providecommand{\AA}{\mathbb{A}}
\providecommand{\BB}{\mathbb{B}}
\providecommand{\CC}{\mathbb{C}}
\providecommand{\DD}{\mathbb{D}}
\providecommand{\EE}{\mathbb{E}}
\providecommand{\FF}{\mathbb{F}}
\providecommand{\GG}{\mathbb{G}}
\providecommand{\HH}{\mathbb{H}}
\providecommand{\NN}{\mathbb{N}}
\providecommand{\PP}{\mathbb{P}}
\providecommand{\QQ}{\mathbb{Q}}
\providecommand{\RR}{\mathbb{R}}
\providecommand{\TT}{\mathbb{T}}
\providecommand{\ZZ}{\mathbb{Z}}
\providecommand{\Gm}{\mathbb{G}_{\mathrm m}}
\providecommand{\Ga}{\mathbb{G}_{\mathrm a}}
\providecommand{\cA}{\mathcal{A}}
\providecommand{\cB}{\mathcal{B}}
\providecommand{\cC}{\mathcal{C}}
\providecommand{\cD}{\mathcal{D}}
\providecommand{\cE}{\mathcal{E}}
\providecommand{\cF}{\mathcal{F}}
\providecommand{\cG}{\mathcal{G}}
\providecommand{\cH}{\mathcal{H}}
\providecommand{\cI}{\mathcal{I}}
\providecommand{\cJ}{\mathcal{J}}
\providecommand{\cK}{\mathcal{K}}
\providecommand{\cL}{\mathcal{L}}
\providecommand{\cM}{\mathcal{M}}
\providecommand{\cN}{\mathcal{N}}
\providecommand{\cO}{\mathcal{O}}
\providecommand{\cP}{\mathcal{P}}
\providecommand{\cQ}{\mathcal{Q}}
\providecommand{\cR}{\mathcal{R}}
\providecommand{\cS}{\mathcal{S}}
\providecommand{\cT}{\mathcal{T}}
\providecommand{\cU}{\mathcal{U}}
\providecommand{\cV}{\mathcal{V}}
\providecommand{\cW}{\mathcal{W}}
\providecommand{\cX}{\mathcal{X}}
\providecommand{\cY}{\mathcal{Y}}
\providecommand{\cZ}{\mathcal{Z}}
\providecommand{\fA}{\mathfrak{A}}
\providecommand{\fB}{\mathfrak{B}}
\providecommand{\fC}{\mathfrak{C}}
\providecommand{\fD}{\mathfrak{D}}
\providecommand{\fE}{\mathfrak{E}}
\providecommand{\fF}{\mathfrak{F}}
\providecommand{\fG}{\mathfrak{G}}
\providecommand{\fH}{\mathfrak{H}}
\providecommand{\fI}{\mathfrak{I}}
\providecommand{\fJ}{\mathfrak{J}}
\providecommand{\fK}{\mathfrak{K}}
\providecommand{\fL}{\mathfrak{L}}
\providecommand{\fM}{\mathfrak{M}}
\providecommand{\fN}{\mathfrak{N}}
\providecommand{\fO}{\mathfrak{O}}
\providecommand{\fP}{\mathfrak{P}}
\providecommand{\fQ}{\mathfrak{Q}}
\providecommand{\fR}{\mathfrak{R}}
\providecommand{\fS}{\mathfrak{S}}
\providecommand{\fT}{\mathfrak{T}}
\providecommand{\fU}{\mathfrak{U}}
\providecommand{\fV}{\mathfrak{V}}
\providecommand{\fW}{\mathfrak{W}}
\providecommand{\fX}{\mathfrak{X}}
\providecommand{\fY}{\mathfrak{Y}}
\providecommand{\fZ}{\mathfrak{Z}}
\providecommand{\fraka}{\mathfrak{a}}
\providecommand{\frakb}{\mathfrak{b}}
\providecommand{\frakc}{\mathfrak{c}}
\providecommand{\frakd}{\mathfrak{d}}
\providecommand{\frakm}{\mathfrak{m}}
\providecommand{\frakp}{\mathfrak{p}}
\providecommand{\frakq}{\mathfrak{q}}
\providecommand{\Hom}{\operatorname{Hom}}
\providecommand{\Ext}{\operatorname{Ext}}
\providecommand{\End}{\operatorname{End}}
\providecommand{\Aut}{\operatorname{Aut}}
\providecommand{\Gal}{\operatorname{Gal}}
\providecommand{\Spec}{\operatorname{Spec}}
\providecommand{\Proj}{\operatorname{Proj}}
\providecommand{\Alg}{\mathbf{Alg}}
\providecommand{\Aff}{\mathbf{Aff}}
\providecommand{\mot}{\mathrm{mot}}
\providecommand{\et}{\mathrm{et}}
\providecommand{\dR}{\mathrm{dR}}
\providecommand{\der}{\mathrm{der}}
\providecommand{\Le}{\leqslant}
\providecommand{\Ge}{\geqslant}
\providecommand{\eps}{\varepsilon}
\providecommand{\mup}{\mu_p}
\makeatletter
\@ifundefined{construction}{\newtheorem{construction}{Construction}}{}
\@ifundefined{notation}{\newtheorem{notation}{Notation}}{}
\@ifundefined{claim}{\newtheorem{claim}{Claim}}{}
\makeatother

% --- Pass-2 extra compile shims ---
\providecommand{\GL}{\operatorname{GL}}
\providecommand{\SL}{\operatorname{SL}}
\providecommand{\PGL}{\operatorname{PGL}}
\providecommand{\Ker}{\operatorname{Ker}}
\providecommand{\Coker}{\operatorname{Coker}}
\providecommand{\rank}{\operatorname{rank}}
\providecommand{\limind}{\varinjlim}
\providecommand{\limproj}{\varprojlim}
\makeatletter
\@ifundefined{scholie}{\newtheorem{scholie}{Scholium}}{}
\makeatother
% --- end pass-2 extra compile shims ---


% --- Pass-2 extra compile shims 2 ---
\providecommand{\SO}{\operatorname{SO}}
\providecommand{\SU}{\operatorname{SU}}
\providecommand{\rk}{\operatorname{rk}}
\providecommand{\fp}{\mathfrak{p}}
\providecommand{\DR}{\mathrm{DR}}
\makeatletter
\@ifundefined{miseengarde}{\newtheorem{miseengarde}{Warning}}{}
\@ifundefined{warning}{\newtheorem{warning}{Warning}}{}
\makeatother
% --- end pass-2 extra compile shims 2 ---


% --- Pass-2 extra compile shims 3 ---
\providecommand{\U}{\operatorname{U}}
\providecommand{\NS}{\operatorname{NS}}
\providecommand{\cris}{\mathrm{cris}}
% --- end pass-2 extra compile shims 3 ---


% --- Pass-2 extra compile shims 4 ---
\providecommand{\na}{a}
\providecommand{\supp}{\operatorname{supp}}
\providecommand{\comp}{\operatorname{comp}}
% --- end pass-2 extra compile shims 4 ---


% --- Pass-2 extra compile shims 5 ---
\providecommand{\dist}{\operatorname{dist}}
\makeatletter
\@ifundefined{variante}{\newtheorem{variante}{Variant}}{}
\makeatother
% --- end pass-2 extra compile shims 5 ---


% --- Pass-2 extra compile shims 6 ---
\providecommand{\R}{\mathrm{R}}
\providecommand{\pr}{\operatorname{pr}}
\makeatletter
\@ifundefined{exemple}{\newtheorem{exemple}{Example}}{}
\makeatother
% --- end pass-2 extra compile shims 6 ---


% --- Pass-2 extra compile shims 7 ---
\providecommand{\codim}{\operatorname{codim}}
\providecommand{\Int}{\operatorname{Int}}
\providecommand{\RHom}{\operatorname{RHom}}
% --- end pass-2 extra compile shims 7 ---


% --- Pass-2 extra compile shims 8 ---
\providecommand{\Lie}{\operatorname{Lie}}
% --- end pass-2 extra compile shims 8 ---


% --- Pass-2 extra compile shims 9 ---
\providecommand{\cl}{\operatorname{cl}}
\providecommand{\Tr}{\operatorname{Tr}}
% --- end pass-2 extra compile shims 9 ---


% --- Pass-2 extra compile shims 10 ---
\providecommand{\comma}{,}
\providecommand{\Pcl}{\operatorname{Pcl}}
\makeatletter
\@ifundefined{paraphrase}{\newtheorem{paraphrase}{Paraphrase}}{}
\makeatother
% --- end pass-2 extra compile shims 10 ---


% --- Pass-2 extra compile shims 11 ---
\providecommand{\re}{\operatorname{Re}}
% --- end pass-2 extra compile shims 11 ---


% --- Pass-2 extra compile shims 12 ---
\providecommand{\ab}{\mathrm{ab}}
\providecommand{\Subset}{\subset}
% --- end pass-2 extra compile shims 12 ---

% --- end pass-2 compatibility shim ---

\begin{document}
\begin{center}\small\emph{Literal English translation working draft. Formulae, numbering, and TeX structure are preserved; this is not yet a modern recast.}\end{center}
\vspace{0.5em}


\title{One the variation under twisting of local constants in functional equations of L-functions}
\author{P. Deligne$^1$ and G. Henniart$^2$\\[0.5em]
\small $^1$Institut des Hautes \'Etudes Scientifiques, 35 route de Chartres, F-91440 Bures-sur-Yvette, France\\
\small $^2$11 rue Ruhmkorff, F-75017 Paris, France}
\date{}
\maketitle

\section*{Introduction}
\addcontentsline{toc}{section}{Introduction}

Soit $K$ un corps local non-archim\'edien, \`a corps r\'esiduel fini. Soient $\overline{K}$ une cl\^oture alg\'ebrique s\'eparable de $K$ et $W = W(\overline{K}/K)$ le groupe de Weil correspondant (cf.~[3], Ap.~II, ou [1], \S 2). La th\'eorie du corps de classes local fournit un isomorphisme $W^{\text{ab}} \xrightarrow{\sim} K^*$, que nous normaliserons comme dans [1], 2.3. Nous utiliserons cet isomorphisme pour identifier syst\'ematiquement classes d'isomorphismes de repr\'esentations de dimension 1 de $W$ et quasi-caract\`eres de $K^*$. Nous fixerons un caract\`ere non trivial $\psi$ du groupe additif de $K$, et une mesure de Haar $dx$ sur ce groupe.

Si $U$ est une repr\'esentation de $W$ (continue et de dimension finie), on sait que pour chaque quasi-caract\`ere $\chi$ de $K^*$ de conducteur $a(\chi)$ suffisamment grand, la constante locale $\varepsilon(U \otimes \chi, \psi, dx)$ (cf.~[1], 4.1) admet une description simple: notant $v_K$ la valuation normalis\'ee de $K$, il existe un \'el\'ement $a$ de $K^*$ tel que l'on ait $\chi(1+y) = \psi(ay)$ si $y \in K$ v\'erifie $2v_K(y) > a(\chi)$, et l'on a, en posant $\gamma = a^{-1}$,
@@MATH0@@

C'est d'ailleurs ce fait qui est \`a la base de la construction des constantes locales donn\'ee dans [1], \S 4. En termes de la repr\'esentation virtuelle $W = U - (\dim U) \cdot 1$, qui est de dimension 0 et de m\^eme d\'eterminant que $U$, la formule \eqref{eq:1} s'\'ecrit
@@MATH1@@

Nous prouvons ici un r\'esultat plus g\'en\'eral. Soit $t$ un nombre r\'eel positif ou nul. Soit $W$ une repr\'esentation virtuelle de $W$, de dimension 0, suppos\'ee triviale sur le groupe de ramification $W(t/2)$ i.e.~dont tous les constituants (cf.~0.5) sont triviaux sur $W(t/2)$ (il s'agit de groupes de ramification en num\'erotation sup\'erieure, cf.~0.4). Soit enfin $V$ une repr\'esentation de $W$, sans vecteur non nul fix\'e par $W(t)$. Alors on a la formule suivante (4.6).
@@MATH2@@
o\`u $\gamma$ est un \'el\'ement de $K^*$ ne d\'ependant que de $V$ et $\psi$. Sa valuation est $a(V) + (\dim V) \cdot n(\psi)$, o\`u $a(V)$ est l'exposant du conducteur de $V$, et $n(\psi)$ l'ordre de $\psi$, cf.~0.3. L'\'el\'ement $\gamma$ n'est bien d\'efini qu'\`a multiplication pr\`es par un \'el\'ement du groupe $U(t/2)$ des unit\'es $u$ de $K$ v\'erifiant $v_K(u-1) > t/2$. Sa classe dans $K^*/U(t/2)$ est d\'ej\`a caract\'eris\'ee par les \'egalit\'es \eqref{eq:3}, quand $W$ parcourt l'ensemble des repr\'esentations de la forme $\eta - 1$, o\`u $\eta$ est un quasi-caract\`ere de $K^*$ trivial sur $U(t/2)$.

Lorsqu'on suppose seulement que $W$ est une repr\'esentation virtuelle de $W$, de dimension 0 et triviale sur $W(t)$, nous prouvons (4.2) que $\varepsilon(W \otimes V, \psi)/\det W(\gamma)$ est une racine de l'unit\'e d'ordre une puissance de $p$.

L'\'el\'ement $\gamma$ de $K^*/U(t/2)$ intervenant dans la formule \eqref{eq:3} se calcule comme suit. Une variante du th\'eor\`eme de Brauer, donn\'ee au \S 2, permet d'\'ecrire $V$, dans le groupe de Grothendieck des repr\'esentations de $W$, comme combinaison lin\'eaire finie de repr\'esentations induites de repr\'esentations $\chi_i$ de dimension 1 de sous-groupes ouverts (d'indice fini) $H_i$ de $W$, la repr\'esentation $\chi_i$ \'etant non-triviale sur $H_i \cap W(t)$:
\[
V = \sum n_i \Ind(\chi_i), \quad n_i \in \mathbb{Z};
\]
a each $H_i$ correspond a extension $L_i$ of $K$, and $\chi_i$ s'identifie a quasi-character of $L_i^*$; one choisit then $a_i$ in $L_i^*$ in such a way that the one ait $\chi_i(1+y) = \psi \circ \Tr_{L_i/K}(a_i y)$ si $y \in L_i$ satisfies $2v_{L_i}(y) > a(\chi_i)$. Si $a(\chi_i) \leq 1$, this condition signifie that $v(a_i) > -n(\psi \circ \Tr_{L_i/K})$, and one impose moreover that $v(a_i) = -(n(\psi \circ \Tr_{L_i/K}) + a(\chi_i))$. Posant $\gamma_i = a_i^{-1}$, one has
\begin{equation}\label{eq:4}
\gamma = \prod N_{L_i/K}(\gamma_i)^{n_i}.
\end{equation}

We ne connaissons not of proof direct of this that the class, in $K^*/U(t/2)$, of the quantite $\gamma = \gamma(V, \psi)$ defined by this formule, ne depend not of the decomposition choisie of $V$ en representations induites.

Si $\chi$ is a representation of dimension 1 of $W$, one can preciser the definition of $\gamma = \gamma(\chi, \psi)$ en exigeant that the one ait
\[
\chi\biggl(\sum_{j=0}^{p-1} \frac{(ay)^j}{j!}\biggr) = \psi(ay)
\]
for $y$ in $K$ verifiant $pv_K(y) > a(\chi)$. One can similarly preciser $\gamma(V, \psi)$, for a representation quelconque $V$ of $W$, en the decomposant as more haut, en definissant a element $\gamma_i$ of $L_i^*$ by the equality
\[
\chi_i\biggl(\sum_{j=0}^{p-1} \frac{(ay)^j}{j!}\biggr) = \psi \circ \Tr_{L_i/K}(a_i y)
\]
for every $y$ in $L_i$ verifiant $pv_{L_i}(y) > a(\chi_i)$, and en definissant $\gamma(V, \psi) = \gamma$ by the formule \eqref{eq:4}. Si $a(\chi_i) \leq 1$, one impose moreover the equality
\[
v(\gamma_i) = n(\psi \circ \Tr_{L_i/K}) + a(\chi_i).
\]
Grace a interpretation en termes of the function $W \mapsto \varepsilon(W \otimes V, \psi, dx)$, we montrons that the element ainsi defined ne depend not of the choix effectues (of the decomposition of $V$ and of the elements $\gamma_i$), a multiplication pres by a element of $U((1-1/p)t)$ (cf.~4.11 a 4.13).

For the formalism auquel $\gamma(V, \psi)$ obeit, see the fin of the chapitre 4 (4.16 and sq.).


\section{Notations}
\label{sec:0}

\setcounter{subsection}{-1}

\subsection{}
\label{sec:0.1}

Si $K$ is a field value complet, one denotes $v$ sa valuation and $\cO$ the ring of valuation corresponding. For every nombre reel $t$, one denotes $A(t)$ the under-group additif of $K$ form of the elements of valuation to the less $t$, and, si $t$ is positive or zero, one denotes $U(t)$ the under-group of the unites $x$ of $\cO$ verifiant $v(x-1) > t$. Si necessary, one precise by a indice $K$. One fixe a nombre prime $p$, and the field values consideres are of caracteristique residuelle $p$.

\subsection{}
\label{sec:0.2}

We shall denote $E$ the exponentielle tronquee:
\[
E(x) = \sum_{i=0}^{p-1} \frac{x^i}{i!},
\]
and $L$ the logarithme tronque
\[
L(x) = \sum_{i=1}^{p-1} (-1)^{i+1} \frac{(x-1)^i}{i}.
\]
For $t$ positive or zero, these maps induisent of the isomorphisms inverses the a of the other:
\[
@@MATH30@@
\]

\subsection{}
\label{sec:0.3}

Let $K$ a field local not-archimedien, i.e.~a field value complet a field residuel finite, fixe in every the sequence. As in [1], we shall denote $\psi: K \to \bC^*$ a character additif not trivial of $K$, $n(\psi)$ the ordre of $\psi$, i.e.~the more grand integer $m$ such that $\psi$ let trivial on $A(-m)$, $dx$ a mesure of Haar on $K$, $\overline{K}$ a algebraic closure separable of $K$ and $W(\overline{K}/K) = W$ the group of Weil (absolute) corresponding.

\subsection{}
\label{sec:0.4}

Si $G$ is the Galois group of a extension galoisienne finite $L/K$, and $t$ a nombre reel positive or zero, one denotes $G(t)$ the group of ramification of indice $t$ en numerotation superieure (cf.~[2], IV, \S 3). Recall that, as function of $t$, $G(t)$ is decroissant, saute en a nombre finite of values, and that between deux sauts successifs $t_1$ and $t_2$, $G(t)$ is constant, egal a $G(t_2)$.

The raison of be of the numerotation superieure is sa compatibility to the passage to the quotient: si $H$ is a under-group distingue of $G$, $(G/H)(t)$ is the image of $G(t)$ in $G/H$. Si $L$ is maintenant a extension galoisienne infinie of $K$ and $G$ the group $\Gal(L/K)$, one defines the groups $G(t)$ by the formule $G(t) = \varprojlim \Gal(E/K)(t)$, or the limite is prise suivant the extensions galoisiennes finite $E$ of $K$ in $L$. The compatibility precedente is conservee. The group of inertie of $G$ n'is other that $G(0)$ and the group of inertie sauvage, son $p$-group of Sylow, is the adherence of the reunion of the $G(t)$ for $t > 0$.

Si $L$ contient a extension not ramifiee maximale of $K$, one defines as in [3], App.~II, the group of Weil $W(L/K)$. It contient the group of inertie of $G$, and this permet of poser $W(L/K)(t) = G(t)$ for $t$ positive or zero.

We shall denote $I$ the group of inertie of $\Gal(\overline{K}/K)$: $I = W(0)$, and $P$ son pro-$p$-under-group of Sylow (group of inertie sauvage).

\subsection{}
\label{sec:0.5}

By representation (of a group topologique), we entendrons toujours representation continue in a space vector complex of dimension finite. In a representation of $W$, the group of inertie $I$ agit a travers a group finite. A representation of $W$ will be dite not ramifiee si $I$ agit trivialement, and ramifiee in the cas contraire. Let $V$ a representation virtuelle of $W$, i.e.~a element of the group of Grothendieck of the category of representations of $W$. It s'ecrit in such a way unique $V = \sum n(V') V'$, or $V'$ parcourt the representations irreductibles of $W$ and $n(V')$ is a integer, the multiplicite of $V'$ in $V$. The $V'$ of multiplicite not zero are en nombre finite; this are the constituants of $V$.

For the definition of the constant local $\varepsilon(V, \psi, dx)$, we renvoyons a [1] 4.1. Recall that for $V$ of dimension 0 (resp.~of dimension 0 and of determinant trivial), it is independante of $dx$ (resp.~of $dx$ and of $\psi$); en this cas we the noterons simplement $\varepsilon(V, \psi)$ (resp.~$\varepsilon(V)$).

One denotes $a(V)$ the exposant of the conducteur of Artin of $V$ and $\sw(V)$ the exposant of son conducteur of Swan. one has
\[
a(V) = \dim V - \dim V^I + \sw(V).
\]

\subsection{}
\label{sec:0.6}

Let $V$ a representation irreducible ramifiee of $W$. Notons $W(V)$ the quotient of $W$ which agit fidelement on $V$; one defines $\alpha(V)$ as the dernier saut of the filtration of $W(V)$: by definition, one has
\[
V^{W(V)_{\alpha(V)-\varepsilon}} = 0 \quad \text{et} \quad V^{W(V)_{\alpha(V)+\varepsilon}} = V
\]
for every $\varepsilon > 0$.

Si $V$ is not ramifiee, one sets $\alpha(V) = 0$. With these conventions, every representation irreducible $V$ of $W$ satisfait a the formule
\[
\sw(V) = \alpha(V) \dim V.
\]

Let $V$ a representation virtuelle of $W$. One denotes $\alpha(V)$ (resp.~$\beta(V)$) the borne inferieure (resp.~superieure) of the $\alpha(V')$, quand $V'$ parcourt the constituants of $V$. That the one ait $\alpha(V) > \alpha$ signifie that the constituants of $V$ n'have not of vecteur not zero fixe by $W(\alpha)$: $V^{W(\alpha)} = 0$. That the one ait $\beta(V) < \beta$ (one has then $\beta > 0$) signifie that $V$ provient by inflation of a representation virtuelle of $W/W(\beta)$.

\subsection{}
\label{sec:0.7}

A finite extension separable $L$ of $K$ defines a set galoisien $\Hom_K(L, \overline{K})$. One sets $\alpha(L/K) = \alpha(U)$, or $U$ is the representation of permutation of $W$ on $\Hom_K(L, \overline{K})$. Si $M$ is a extension galoisienne of $K$ contenant $L$, $\alpha(L/K)$ is the borne inferieure of the indices $\alpha$ such that $\Gal(M/L)$ contienne $\Gal(M/K)(\alpha)$.

One sait (cf.~[2], XV) that the isomorphism of the theorie of the field of classes local: $W^{\text{ab}} \xrightarrow{\sim} K^*$ identifie $W^{\text{ab}}(t)$ a $U(t)$. Si $\chi$ is a quasi-character ramifie of $K^*$ (corresponding a representation of $W$ of dimension 1), $\alpha(\chi)$ is therefore the more grand integer $n > 0$ such that $\chi$ let not trivial on $U_K(n)$. Si $\chi$ is a quasi-character of $K^*$, not ramifie or moderement ramifie, one has $\alpha(\chi) = 0$. One dit that $\chi$ is sauvagement ramifie si the one has $\alpha(\chi) > 0$.

\subsection{}
\label{sec:0.8}

One denotes $\bQ_p/\bZ_p(1)$ the under-group of $\overline{\bQ}_\ell^*$ form of the racines of the unite of ordre a power of $p$. One denotes $\mod^*$ a congruence multiplicative, i.e.~a congruence in the group multiplicatif of a field. Si $X$ is a set finite, one denotes $|X|$ son cardinal.


\section{The cas of dimension 1}
\label{sec:1}
\setcounter{equation}{0}

\begin{proposition}\label{prop:1.1}
Let $\chi$ and $\eta$ deux quasi-characters of $K^*$, verifiant $\alpha(\eta) < \alpha(\chi)$. Let $m$ the more petit integer such that the one ait $2m > \alpha(\chi)$ and let $a$ a element of $K$ such that the one ait $\chi(1+y) = \psi(ay)$ of the that the element $y$ of $K$ is of valuation to the less $m$. one has then $v_K(a) = -(n(\psi) + 1 + \alpha(\chi))$ and $a$ is unique $\mod^*\, U(\alpha(\chi)/2)$. Similarly, let $b$ a element of $K$ verifiant $\eta(1+y) = \psi(by)$ of the that $v(y) \geq m$. Then one has
\begin{equation}\label{eq:1.1.1}
\varepsilon((\eta-1)\chi, \psi) = \eta^{-1}(a) \cdot (\chi\eta)^{-1}(1+b/a) \cdot \psi(b).
\end{equation}
\end{proposition}

\begin{proof}
Remplacer $\chi$ by $\chi\eta$ ne change not the conducteur, and mene a remplacer $a$ by $a+b$, similarly valuation. Si $\alpha(\chi)$ is pair, one has therefore:
\[
\frac{\varepsilon(\eta\chi, \psi, dx)}{\varepsilon(\chi, \psi, dx)} = \frac{(\chi\eta)^{-1}(a+b) \, \psi(a+b)}{\chi^{-1}(a) \, \psi(a)},
\]
of or
\[
\varepsilon((\eta-1)\chi, \psi) = \eta^{-1}(a) \cdot (\chi\eta)^{-1}(1+b/a) \cdot \psi(b).
\]
For montrer that the same result subsiste si $\alpha(\chi)$ is impair, it is enough to verify that en this cas the rapport
\[
R(u) = (\chi\eta)^{-1}((a+b)u) \, \psi((a+b)u) \big/ \chi^{-1}(au) \, \psi(au)
\]
is independant of $u$ variant in $U(m-1)$. Or one has
\[
R(u)/R(1) = \eta^{-1}(u) \, \psi(b(u-1)),
\]
and, puisque $\alpha(\chi)$ is impair, the inegalite $2m > \alpha(\chi)$ entraine $(m-1) \geq \alpha(\chi)/2$, and the definition of $b$ assure that $R(u)/R(1)$ vaut 1. C.Q.F.D.
\end{proof}

\begin{corollary}\label{cor:1.3}
Conservons the hypotheses and notations of \ref{prop:1.1}.
\begin{enumerate}[label=(\roman*)]
\item one has $\varepsilon((\eta-1)\chi, \psi) \equiv \psi(a^{-1}) \pmod{*} \;\bQ_p/\bZ_p(1)$.
\item Si moreover one has $2\alpha(\eta) < \alpha(\chi)$, one has $\varepsilon((\eta-1)\chi, \psi) = \psi(a^{-1})$.
\end{enumerate}
\end{corollary}

Puisque the one has $\alpha(\eta) < \alpha(\chi)$, one has also $v(b) > v(a)$ and $1+b/a$ appartient a $U(1)$. One proves (i) en notant that $\psi$, and the restriction a $U(1)$ of a quasi-character quelconque of $K^*$, prennent values in $\bQ_p/\bZ_p(1)$. The condition $\alpha(\eta) < \alpha(\chi)/2$ permet of prendre $b = 0$ in \ref{prop:1.1}, this which proves (ii).

\begin{corollary}\label{cor:1.4}
Let $\chi$ a quasi-character sauvagement ramifie of $K^*$, $m$ the more petit integer such that the one ait $2m > \alpha(\chi)$, and $a$ a element of $K^*$ such that the one ait $\chi(1+y) = \psi(ay)$ si $y \in K$ is of valuation to the less $m$. Let $(\eta_i)_{i \in I}$ a family finite, of to the less deux elements, of quasi-characters of $K^*$, verifiant $\alpha(\eta_i) < \alpha(\chi)$. One choisit of the elements $b_i$ of $K$ such that $\eta_i(1+y) = \psi(b_i y)$ quand $v(y) \geq \alpha(\chi)/2$. For each part $J$ of $I$, one sets
\[
\eta(J) = \prod_{j \in J} \eta_j, \quad b(J) = \sum_{j \in J} b_j, \quad \varepsilon(J) = (-1)^{|J|}.
\]
Then one has
\begin{equation}\label{eq:1.4.1}
\varepsilon\biggl(\chi \prod_{i \in I}(1-\eta_i)\biggr) = \eta(I)^{-1}\biggl(\prod_{J \subset I} (1+b(J)/a)^{\varepsilon(J)}\biggr) \cdot \prod_{i \in I} \eta_i^{-1}\biggl(\prod_{J \subset I \setminus \{i\}} (1+b(J)/a)^{\varepsilon(J)}\biggr).
\end{equation}
\end{corollary}

\begin{remark}
Of the that $I$ a to the less deux elements, the representation virtuelle $\chi \prod_{i \in I}(1-\eta_i)$, of dimension 0, is of determinant trivial. C'is this which a permis of omettre $\psi$ under the signe $\varepsilon$. of ailleurs, in the membre of line, outre $\chi$ and the $\eta_i$, seuls apparaissent the quotients $b_i/a$, eux also independants of $\psi$.
\end{remark}

\subsection{}
\label{sec:1.5}

Demontrons the corollary \ref{cor:1.4}. one has
\[
\prod_{i \in I}(1-\eta_i) = \sum_{J \subset I} \varepsilon(J) \eta(J) = \sum_{J \subset I} \varepsilon(J)(\eta(J)-1),
\]
of or
\[
\varepsilon\biggl(\chi \prod_{i \in I}(1-\eta_i)\biggr) = \prod_{J \subset I} \varepsilon((\eta(J)-1)\chi, \psi)^{\varepsilon(J)}.
\]
Each factor is justifiable of \eqref{eq:1.1.1}, and the one can choisir $b(J)$ as element $b$ attached to $\eta(J)$. Puisque $\prod_{J \subset I} \eta(J)^{\varepsilon(J)}$ vaut 1 and that $\sum_{J \subset I} \varepsilon(J) b(J)$ vaut 0, one trouve
\[
\varepsilon\biggl(\chi \prod_{i \in I}(1-\eta_i)\biggr) = \prod_{J \subset I} (\chi\eta(J))^{-1}(1+b(J)/a) \, \psi(b(J))^{\varepsilon(J)}.
\]
One can ecrire the membre of line as the product of
\[
\chi^{-1}\biggl(\prod_{J \subset I} (1+b(J)/a)^{\varepsilon(J)}\biggr)
\]
and of the product on $i$ of the quantites
\[
\eta_i^{-1}\biggl(\prod_{J \subset I} (1+b(J)/a)^{\varepsilon(J)}\biggr),
\]
egales a
\[
\eta_i^{-1}\biggl(\prod_{J \subset I \setminus \{i\}} (1+b(J)/a)^{\varepsilon(J)} \cdot \prod_{J \subset I \setminus \{i\}} (1+b(J)/a)^{-\varepsilon(J)}\biggr).
\]
That proves the formule \eqref{eq:1.4.1}.

\begin{lemma}\label{lem:1.6}
Let $I$ a set finite not vide. Let $t = (t_i)_{i \in I}$ a family of indeterminees and $\bZ[t]$ the ring of the series formelles en the $t_i$, a coefficients integers. For each under-set $J$ of $I$, one sets $t(J) = \sum_{j \in J} t_j$ and $\varepsilon(J) = (-1)^{|J|}$. In $\bZ[t]$, the product $\prod_{J \subset I} (1+t(J))^{\varepsilon(J)}$ is congru a $1$ modulo $\prod_{i \in I} t_i$.
\end{lemma}

This assertion resulte of this that the product considered devient identiquement 1 si one fait $t_i = 0$ for a indice $i$.

\begin{remark}
The lemma \ref{lem:1.11} precisera the lemma \ref{lem:1.6}.
\end{remark}

\begin{proposition}\label{prop:1.7}
Let $\chi$ a quasi-character sauvagement ramifie of $K^*$. Let $(\eta_i)_{i \in I}$ a family finite of quasi-characters of $K^*$. One assumed that the one has $\alpha(\eta_i) < \alpha(\chi)$ for each indice $i$ and $\sum_{i \in I}(\alpha(\chi) - \alpha(\eta_i)) > \alpha(\chi)$. Then one has
\[
\varepsilon\biggl(\chi \prod_{i \in I}(1-\eta_i)\biggr) = 1.
\]
\end{proposition}

\begin{proof}
the hypothese assure that $I$ a to the less deux elements. One can then appliquer the corollary \ref{cor:1.4}. By the lemma \ref{lem:1.6}, one has
\[
\prod_{J \subset I} (1+b(J)/a)^{\varepsilon(J)} \equiv 1 \pmod{*} \prod_{i \in I} (b_i/a).
\]
But one can choisir the $b_i$ in such a way that the valuation of $b_i/a$ let $\alpha(\chi) - \alpha(\eta_i)$ (cf.~\ref{prop:1.1}). Celle of the product of the $b_i/a$ vaut then to the less $\alpha(\chi)+1$ and $\chi\eta(J)$ prend the value 1 on $\prod_{J \subset I}(1+b(J)/a)^{\varepsilon(J)}$. Similarly, for each indice $i$, one has
\[
\prod_{J \subset I \setminus \{i\}} (1+b(J)/a)^{\varepsilon(J)} \equiv 1 \pmod{*} \prod_{j \in I \setminus \{i\}} (b_j/a),
\]
and the valuation of this dernier product vaut to the less $\alpha(\chi) + 1$. one has therefore $\eta_i^{-1}\bigl(\prod_{J \subset I \setminus \{i\}}(1+b(J)/a)^{\varepsilon(J)}\bigr) = 1$ and \eqref{eq:1.4.1} permet of conclure.
\end{proof}

Recall the

\begin{definition}\label{def:1.8}
Let $G$ and $H$ deux groups abeliens, notes additivement. A function $f: G \to H$ is dite of degree to the more $n$ si for every family $(x_i)_{i \in I}$ of $n+1$ elements of $G$, one has
\begin{equation}\label{eq:1.8.1}
\sum_{J \subset I} (-1)^{|J|} f\biggl(\sum_{j \in J} x_j\biggr) = 0.
\end{equation}
We shall say that $f$ is homogene of degree $n$ si $f$ is of degree to the more $n$ and that en outre one has $f(kx) = k^n f(x)$ for every integer $k$.
\end{definition}

Let $x \mapsto \delta(x)$ the map canonical of $G$ in the algebra of group $\bZ[G]$. Notons $\tilde{f}$ the homomorphisme of $\bZ[G]$ in $H$ defined by $f$: $\tilde{f} \circ \delta = f$. The condition \eqref{eq:1.8.1} s'ecrit again
\begin{equation}\label{eq:1.8.2}
\tilde{f}\biggl(\prod_{i \in I}(1-\delta(x_i))\biggr) = 0.
\end{equation}
In $\bZ[G]$, one has
\begin{multline*}
(1-\delta(u))(1-\delta(x))(1-\delta(y)) = (1-\delta(x))(1-\delta(y)) - (1-\delta(u))(1-\delta(y)) \\
- (1-\delta(y))(1-\delta(x)) + (1-\delta(x+y))(1-\delta(u)),
\end{multline*}
of sorte that \eqref{eq:1.8.2} signifie also that $\tilde{f}$ s'annule on the power $(n+1)$\textsuperscript{eme} of the ideal of augmentation of $\bZ[G]$. In particular, si $f$ is of degree to the more $n$, it is of degree to the more $m$ for every $m$ superieur a $n$.

\begin{lemma}\label{lem:1.9}
Let $G$ a group abelien, $H$ a $\bZ[1/n!]$-module, and $f$ a function of $G$ in $H$, of degree to the more $n$. One can, in such a way unique, decomposer $f$ en a sum $\sum_{i=0}^n f_i$, or $f_i$ is homogene of degree $i$.
\end{lemma}

It is easy of see that $f$ admet to the more a telle decomposition. Indeed, the systeme of equations lineaires
\[
\sum_{i=0}^n k^i f_i(x) = f(kx), \quad k \text{ variant de } 0 \text{ a } n,
\]
determine the $f_i(x)$, puisque the determinant of Vandermonde $\det(k^i)$ is inversible in $\bZ[1/n!]$.

Prouvons the existence of the decomposition, by recurrence on $n$. For $n=0$, $f$ is a function constant and $f = f_0$ is the decomposition cherchee. Supposons $n$ not zero, and let
\[
f(x_1, \ldots, x_n) = \tilde{f}\bigl((\delta(x_1)-1) \cdots (\delta(x_n)-1)\bigr)
\]
the $n$\textsuperscript{eme} polarise of $f$.

This function vaut 0 of the that a of the $x_i$ s'annule, and is lineaire en each $x_i$; indeed, one has
\begin{align*}
f(x+y, x_2, \ldots, x_n) &= \tilde{f}\bigl((\delta(x+y)-1)(\delta(x_2)-1)\cdots(\delta(x_n)-1)\bigr) \\
&= \tilde{f}\bigl((\delta(x)-1)(\delta(x_2)-1)\cdots(\delta(x_n)-1)\bigr) + \tilde{f}\bigl((\delta(y)-1)(\delta(x_2)-1)\cdots(\delta(x_n)-1)\bigr) \\
&= f(x, x_2, \ldots, x_n) + f(y, x_2, \ldots, x_n).
\end{align*}
Posons $f_n(x) = f(x, \ldots, x)$. This function is homogene of degree $n$, of $n$\textsuperscript{eme} polarisee $f(x_1, \ldots, x_n)$. The function $f - f_n$ is of degree to the more $n-1$, and one applique the hypothese of recurrence.

\subsection{}
\label{sec:1.10}

For $\chi$ fixe, \ref{prop:1.7} shows that the function $\eta \mapsto \varepsilon(\eta\chi, \psi, dx)$, of the group of the quasi-characters of $K^*/U((1-1/n)\alpha(\chi))$ in $\bC^*$, is of degree $< n$. The proposition which suit explicite a decomposition \ref{lem:1.9} lorsque $n \leq p$.

\begin{proposition}\label{prop:1.10}
Let $\chi$ and $\eta$ deux quasi-characters of $K^*$. Let $n$ a integer, $2 \leq n \leq p$. One assumed that $\chi$ is sauvagement ramifie and that $\eta$ satisfies $\alpha(\eta) < (1-1/n)\alpha(\chi)$. Let $a$ and $b$ deux elements of $K$ such that the one ait $\chi(Ey) = \psi(ay)$ for $v(y) > \alpha(\chi)/n$ and $\eta(Ey) = \psi(by)$ for $v(y) \geq \alpha(\chi)/n$. one has
\begin{equation}\label{eq:1.10.1}
\varepsilon(\eta\chi, \psi) = \varepsilon(\chi, \psi) \cdot \eta^{-1}(a) \cdot \psi\biggl(\sum_{i=1}^{n-1} \frac{(-1)^{i-1}}{i} \, b^i/a^{i-1}\biggr).
\end{equation}
\end{proposition}

Si $\eta$ is not trivial on $U(\alpha(\chi)/n)$, one has $\alpha(\chi) - \alpha(\eta) = v(b/a)$. Otherwise, one can only conclure that the one has $v(b) > -n(\psi) - \alpha(\chi)/n$, of or $v(b/a) > (1-1/n)\alpha(\chi)$. In all the cas, the hypothese $\alpha(\eta) < (1-1/n)\alpha(\chi)$, or again $n(\alpha(\chi) - \alpha(\eta)) > \alpha(\chi)$, entraine $nv(b/a) > \alpha(\chi)$; it assure that $\eta$ is trivial on $U(v(b^n/a^n))$ and $\psi$ on $A(v(b^n/a^{n-1}))$. Puisque $n$ vaut to the more $p$, one has in particular
\[
1 + b/a = L(1+b/a) = E\biggl(\sum_{i=1}^{p-1} \frac{(-1)^{i-1}}{i} (b/a)^i\biggr) \pmod{*} U(\alpha(\chi)+1),
\]
of or
\[
\eta^{-1}(a)\eta^{-1}(1+b/a) = \eta^{-1}\bigl((a+b) \cdot L(1+b/a)\bigr) \cdot \psi(b).
\]
One calcule aisement
\begin{align*}
(a+b) \sum_{i=1}^{p-1} \frac{(-1)^{i-1}}{i} (b/a)^i &= -b + \sum_{i=2}^{p-1} (-1)^{i-1}\Bigl(\frac{1}{i} - 1\Bigr) \frac{b^i}{a^{i-1}} \\
&\quad + (-1)^{p-1}\frac{1}{p}\binom{p}{p-1}\frac{b^p}{a^{p-1}}.
\end{align*}
As the dernier terme, and all ceux of indice $i \geq n$, are negligeables under the signe $\psi$, the formule \eqref{eq:1.10.1} resulte of \eqref{eq:1.1.1}.

\subsection{}
\label{sec:1.11}

The fin of this n\textsuperscript{o} ne servira more in what follows of the article. We y donnons a generalisation commune of \ref{prop:1.7} and \ref{prop:1.10}.

\begin{lemma}\label{lem:1.11}
One conserve the notations of the lemma \ref{lem:1.6}. For each family finite of integers natural $n = (n_i)_{i \in I}$, one posera $|n| = \sum_{i \in I} n_i$, $n! = \prod_{i \in I} n_i!$ and $t^n = \prod_{i \in I} t_i^{n_i}$; one notera $J(n)$ the set of the indices $i$ such that $n_i$ let not zero. One appelle $N$ the cardinal of $I$.
\begin{enumerate}[label=(\roman*)]
\item In $\bZ[t]$, $\prod_{i \in I}(1+t_i)$ is sum of $(-1)^{N-1}(N-1)! \prod_{i \in I} t_i$ and of termes of degree (total) to the less $N+1$.
\item In $\bZ_{(p)}[t]$, one has
\begin{equation}\label{eq:1.11.1}
\prod_{J \subset I}(1+t(J))^{\varepsilon(J)} \equiv \sum_n \frac{(-1)^{|n|-1}(|n|-1)!}{n!} \, t^n \pmod{\prod_{i \in I} t_i^p},
\end{equation}
the sum portant on the families of integers natural $n$ such that $n_i > 0$ for every $i$ and that to the less a of the $n_i$ let strictement inferieur a $p$.
\end{enumerate}
\end{lemma}

Demontrons this lemma. It is enough to prouver (i) in $\bQ[t]$, or one has
\[
\log \prod_{J \subset I}(1+t(J))^{\varepsilon(J)} = \sum_{J \subset I} \varepsilon(J) \sum_{n>0} \frac{(-1)^{n-1}}{n} t(J)^n = \sum_{n>0} \frac{(-1)^{n-1}}{n} \sum_{J \subset I} \varepsilon(J) t(J)^n.
\]
Developpons $t(J)^n$ selon the formule of the multinome. Then the coefficient of the monome $t^n$ (or $|n| = n$) in $\sum_{J \subset I} \varepsilon(J) t(J)^n$ vaut $\sum_{J \subset I} \varepsilon(J) \, n!/n!$, zero si $J(n)$ is distinct of $I$.
one has therefore
\[
\sum_{J \subset I} \varepsilon(J) t(J)^n = (-1)^N \sum_{\substack{|n|=n \\ J(n)=I}} \frac{n!}{n!},
\]
or the sum porte on $N$-uples $n$ verifiant $|n| = n$ and $J(n) = I$. One en tire
\begin{equation}\label{eq:1.11.2}
\log \prod_{J \subset I}(1+t(J))^{\varepsilon(J)} = \sum_n (-1)^{|n|} \frac{(|n|-1)!}{n!} \, t^n,
\end{equation}
and this expression is bien sum of $-(N-1)! \prod_{i \in I} t_i$ and of termes of degrees superieurs. Prenant the exponentielle, one obtains (i), and a new proof of the lemma \ref{lem:1.6}.

For demontrer (ii), one denotes all of abord that the deux membres of \eqref{eq:1.11.1} are bien with coefficients in $\bZ_{(p)}$; a line, one has indeed, for each indice $i$,
\[
\frac{(|n|-1)!}{n!} = \frac{1}{n_i} \cdot \frac{(|n|-1)!}{(n_i-1)! \prod_{j \neq i} n_j!} \in \bZ
\]
si one choisit $i$ of sorte that $n_i$ let to the more $p-1$, one obtains $(|n|-1)!/n! \in \bZ_{(p)}$. It is enough therefore a new of verify the congruence in $\bQ[[t]]$. Grace to the lemma \ref{lem:1.6}, one has, in $\bQ[[t]]$,
\[
\prod_{J \subset I}(1+t(J))^{\varepsilon(J)} \equiv \sum_n \frac{(-1)^{|n|-1}(|n|-1)!}{n!} \, t^n \pmod{\prod_{i \in I} t_i^p}.
\]
One utilise then \eqref{eq:1.11.2}. C.Q.F.D.

\begin{proposition}\label{prop:1.11}
Let $\chi$ a quasi-character sauvagement ramifie of $K^*$. One fixe a element $a$ of $K^*$ such that the one ait $\chi(Ey) = \psi(ay)$ of the that $y \in K$ satisfies $pv(y) > \alpha(\chi)$. One se given a family finite $(\eta_i)_{i \in I}$ of quasi-characters of $K^*$, and a family finite $(\alpha_i)_{i \in I}$ of nombres reels. One assumed that for each indice $i$, one has
\[
(\alpha(\chi) - \alpha_i) + p \sum_{j \in I \setminus \{i\}} (\alpha(\chi) - \alpha_j) > \alpha(\chi)
\]
and
\[
\sum_{j \in I \setminus \{i\}} (\alpha(\chi) - \alpha_j) > \alpha_i.
\]
One fixe of the elements $b_i$ of $K$ such that the one ait $\eta_i(Ey) = \psi(b_i y)$ of the that the one has $pv(y) > \alpha_i$. With the notations of the lemma \ref{lem:1.11}, one has then
\begin{equation}\label{eq:1.11.3}
\varepsilon\biggl(\chi \prod_{i \in I}(1-\eta_i)\biggr) = \psi\biggl(\sum_n \frac{(-1)^{|n|-1}(|n|-1)!}{n!} \, b^n/a^{|n|-1}\biggr),
\end{equation}
the sum portant on the families $n = (n_i)_{i \in I}$ such that each $n_i$ let $> 0$ and that to the less deux of the $n_i$ let strictement inferieurs a $p$.
\end{proposition}

\begin{remark}
This proposition permet of donner a other proof of the proposition \ref{prop:1.7}. Under the hypotheses of \ref{prop:1.7} and with the same notations, one can poser $\alpha_i = \alpha(\eta_i)$ and appliquer \eqref{eq:1.11.3}. The termes of the sum of line are of valuation to the less celle of $a^{-1} \prod_{i \in I}(b_i/a)$. The $b_i$ verifiant $v(b_i/a) \geq \alpha(\chi) - \alpha(\eta_i)$, $\chi$ is trivial on $U(v(\prod_{i \in I}(b_i/a)))$, therefore $\psi$ the is on $A(v(a^{-1}\prod_{i \in I}(b_i/a)))$. The membre of line of \eqref{eq:1.11.3} is therefore trivial.
\end{remark}

\subsection{}
\label{sec:1.12}

Demontrons the proposition \ref{prop:1.11}. The hypotheses entrainent that there is to the less deux characters $\eta_i$. one has $v(b_i/a) \geq \alpha(\chi) - \alpha_i$, and of after the lemma \ref{lem:1.6}
\[
\prod_{J \subset I}(1+b(J)/a)^{\varepsilon(J)} \in U\biggl(\sum_{i \in I}(\alpha(\chi) - \alpha_i)\biggr),
\]
\[
\prod_{J \subset I \setminus \{i\}}(1+b(J)/a)^{\varepsilon(J)} \in U\biggl(\sum_{j \in I \setminus \{i\}}(\alpha(\chi) - \alpha_j)\biggr).
\]
For every part $I_1$ of $J$, posons
\[
T(I_1) = L\biggl(\prod_{J \subset I_1}(1+b(J)/a)^{\varepsilon(J)}\biggr).
\]
By hypothese, one has
\[
p \sum_{j \in I \setminus \{i\}}(\alpha(\chi) - \alpha_j) > \alpha(\chi)
\]
and
\[
p \sum_{j \in I}(\alpha(\chi) - \alpha_j) > \alpha_i;
\]
one can therefore ecrire the formule \eqref{eq:1.4.1} under the form
\[
\varepsilon\biggl(\chi \prod_{i \in I}(1-\eta_i)\biggr) = \psi\bigl(-(aT + \sum_i b_i T(I \setminus \{i\}))\bigr).
\]
As one has $p \sum_{j \in I}(\alpha(\chi) - \alpha_j) > \alpha_i$, one can, under the signe $\psi$, negliger for each $i$ the multiples of
\[
b_i \prod_{j \in I \setminus \{i\}}(b_j/a)^{n_j},
\]
i.e.~negliger the multiples of the monomes $b^n/a^{|n|-1}$ lorsque $J(n) = I$ and that to the more a of the $n_i$ satisfies $n_i < p$.

One can therefore utiliser the congruence \eqref{eq:1.11.1} for remplacer $T(J)$, under the signe $\psi$, by
\[
T = \sum_n (-1)^{|n|-1} \frac{(|n|-1)!}{n!} \, b^n/a^{|n|},
\]
or the sum porte on the $n$ such that $J(n) = I$, and that to the less deux of the $n_i$ let strictement inferieurs a $p$.

Similarly, one can remplacer $T(I \setminus \{i\})$ by
\[
S_i = \sum_n (-1)^{|n|-1} \frac{(|n|-1)!}{n!} \, b^n/a^{|n|},
\]
or the sum porte on the $n$ such that $J(n) = I \setminus \{i\}$, and that to the less deux of the $n_j$ let strictement inferieurs a $p$.

Regroupant $T$ and $S_i$, and changeant $n_i$ en $n_i + 1$, one sees that one can remplacer $b_i(T + S_i)$ by
\[
\sum_n (-1)^{|n|-2} \frac{(|n|-2)!}{n!} \, b^n/a^{|n|-1},
\]
or the sum porte on the $n$ such that $J(n) = I$, and that deux to the less of the $n_i$ satisfy $n_i < p$.

Regroupant enfin with the terme $aT$, one has remplace the expression under the signe $\psi$ by
\[
\sum_n (-1)^{|n|-1} \frac{(|n|-1)!}{n!} \, (|n|-1) \, b^n/a^{|n|-1},
\]
and the one sum on the same $n$. As one has $|n| = \sum_{i \in I} n_i$, \eqref{eq:1.11.3} en resulte.


\section{A variante of the theorem of Brauer}
\label{sec:2}

\subsection{}
\label{sec:2.1}

Let $G$ a group finite. We shall denote $R(G)$ the group of Grothendieck of the category of representations of $G$. Si $H$ is a under-group of $G$, the restriction a $H$ of a representation of $G$, and induction a $G$ of a representation of $H$ fournissent of the homomorphismes
\[
\Res_H^G: R(G) \to R(H) \quad \text{et} \quad \Ind_H^G: R(H) \to R(G).
\]

Si $H$ is distingue in $G$, the action of $G$ on $H$ by automorphismes interieurs fournit a action of $G$ on the set $\hat{H}$ of the classes of isomorphism of representations irreductibles of $H$. Si $\tau$ is a telle class, we shall denote $Z(\tau)$ the fixateur of $\tau$ for this action; c'is a under-group of $G$ contenant $H$. For each representation $V$ of $G$, the composante $\tau$-isotypique $V_\tau$ of $V$ is stable by $Z(\tau)$. The functor $V \mapsto V_\tau$ fournit a homomorphisme $V \mapsto v_\tau$ of $R(G)$ in $R(Z(\tau))$, and one checks that the one has
\begin{equation}\label{eq:2.1.1}
v = \sum_{\tau \in \hat{H}/G} \Ind_{Z(\tau)}^G(v_\tau).
\end{equation}

\begin{proposition}\label{prop:2.2}
\begin{enumerate}[label=(\roman*)]
\item Let $G$ a group finite, $A$ a under-group abelien distingue of $G$, and $V$ a representation of $G$. There exists of the under-groups $H_i$ of $G$, contenant $A$, of the characters $\chi_i$ of the $H_i$ and of the integers $n_i$ such that, in $R(G)$, one ait
\[
V = \sum n_i \Ind_{H_i}^G(\chi_i).
\]
\item Si the space $V^A$ of the points of $V$ fixes by $A$ is trivial, one can choisir the decomposition of sorte that $\chi_i$ let not trivial on $A$.
\end{enumerate}
\end{proposition}

\begin{proof}
Supposons of abord $A$ central in $G$. One can supposer $V$ irreducible, auquel cas $A$ agit on $V$ by a character $\xi$. The theorem of Brauer assure the existence of decompositions
\begin{equation}\label{eq:2.2.1}
V = \sum n_i \Ind_{H_i}^G(\chi_i), \quad n_i \in \bZ.
\end{equation}
Ecrivons that the one has $V = \sum V_\tau$:
\[
V = \sum n_i \Ind_{H_i}^G(\chi_i) = \sum n_i \Ind_{H_i A}^G\bigl((\Ind_{H_i}^{H_i A} \chi_i)_\xi\bigr).
\]
The $\xi$-composante isotypique of $\Ind_{H_i}^{H_i A} \chi_i$ is trivial si $\chi_i$ is distinct of $\xi$ on $H_i \cap A$, and is the unique character of $H_i A$ prolongeant $\chi_i$ and $\xi$, s'they co\"incident on $H_i \cap A$. This we given a decomposition \eqref{eq:2.2.1}, or the $H_i$ contiennent $A$ and the $\chi_i$ prolongent $\xi$. one has therefore (i) and (ii).

In the cas general, ecrivons the formule \eqref{eq:2.1.1}
\[
v = \sum_{\xi \in \hat{A}/G} \Ind_{Z(\xi)}^G(v_\xi).
\]
Under the hypothese (ii), one can omettre the character $\xi = 0$ of the sum, the space $V_0$ corresponding being zero. On $V_\xi$, $A$ agit selon $\xi$. The representation of $Z(\xi)$ on $V_\xi$ se factorise by $Z(\xi)/\Ker(\xi)$, and in this group $A/\Ker(\xi)$ is central; this we ramene to the cas central, deja traite.
\end{proof}

\begin{proposition}[Variante]\label{prop:2.3}
Let $G$ a group, extension of $\bZ$ by a group finite $G(0)$. Let $A$ a under-group abelien distingue of $G$, contenu in $G(0)$, and $V$ a representation of $G$. There exists of the under-groups of indice finite $H_i$ of $G$ (en nombre finite), contenant $A$, of the characters $\chi_i$ of the $H_i$ and of the integers $n_i$ such that, in $R(G)$, one ait $V = \sum n_i \Ind_{H_i}^G(\chi_i)$. Si $V^A$ is zero, one can choisir the decomposition in such a way that $\chi_i$ let not trivial on $A$.
\end{proposition}

One can supposer $V$ irreducible. One verifies facilement (cf.~[1] 4.10) that every representation irreducible of $G$ is the tensor product of a representation which se factorise by a quotient finite $\bar{G}$ of $G$, by a representation of dimension 1 trivial on $G(0)$. It is enough to appliquer \ref{prop:2.2} a $\bar{G}$.


\section{Normes and traces}
\label{sec:3}

\subsection{}
\label{sec:3.1}

For each finite extension separable $L$ of $K$, notons $J_L$ the group $L^*/\cO_L^*$; the valuation $v_L$ the identifie a $\bZ$. Let $T_L$ the demi-line of the elements positifs or nuls of $J_L \otimes_\bZ \bR$; the valuation $v_L$ the identifie a $\bR^+$.

Recall the methode signalee by Serre ([2], IV \S 3, Rem.~3 p.~83) for indexer the groups of ramification: the functions $\phi$ and $\psi$ of Herbrand engendrent a systeme transitif of isomorphisms between the $T_L$. Appelons $T$ the projective limit of the $T_L$, and notons again $v_L$ the isomorphism compose $T \to T_L \to \bR^+$. Si $L$ is a finite extension separable of $K$, and $M$ a extension galoisienne of $L$, the groups of ramification of $\Gal(M/L)$ are naturally indexes by $T$; one denotes $\Gal(M/L)[t]$ the group of indice $t$. This indexation is compatible to the passages to the under-groups and to the groups quotients. One retrouve the numerotation superieure en identifiant $T$ a $\bR^+$ by $v_K$. Si $M$ is a finite extension of $L$, one retrouve the numerotation inferieure en identifiant $T$ a $\bR^+$ by $v_M$; one passe then of the numerotation superieure a the numerotation inferieure by the function of Herbrand $\Psi_{M/L} = v_M \circ \psi_{M/L} \circ v_L^{-1}$. This function is a homeomorphisme of $\bR^+$ on lui-same, lineaire by morceaux and convexe; sa derivee a gauche en $v_L(t)$, or $t$ appartient a $T$ is the indice of $\Gal(M/L)[t]$ in $\Gal(M/L)[0]$.

Quelles that let the extensions separables finite $L$ and $M$ of $K$, we poserons $\Psi_{M/L} = v_M \circ \psi_{M/L} \circ v_L^{-1}$. Si $M$ is a extension of $L$, we poserons $\psi_{M/L} = \Psi_{M/L}^{-1}$.

\begin{proposition}\label{prop:3.2}
Let $L$ a finite extension separable of $K$, $M$ a finite extension separable of $L$.
\begin{enumerate}[label=(\roman*)]
\item The function $\Psi_{M/L}$ is convexe.
\item Si the extension $M/L$ is moderee (i.e.~si $\alpha(M/L)$ is zero cf.~0.6), $\Psi_{M/L}$ is lineaire. Otherwise, the dernier saut of the derivee of $\Psi_{M/L}$ se product en $\alpha(M/L)$.
\item Let $t \in \bR^+$. Si $t > \alpha(M/L)$, the pente of $\Psi_{M/L}$ en $t$ is the indice of ramification $e$ of $M/L$. Si $t < \alpha(M/L)$, c'is a integer of the form $e/p^k$, $k \geq 1$. En consequence, one has
\[
\Psi_{M/L}(\alpha(M/L)) \leq (e/p) \, \alpha(M/L).
\]
\end{enumerate}
\end{proposition}

\begin{proof}
Let $N$ a extension galoisienne finite of $L$ contenant $M$; posons $G = \Gal(N/L)$ and $H = \Gal(N/M)$. By definition, one has $\Psi_{M/L} = \Psi_{N/L} \circ \Psi_{N/M}^{-1}$, of sorte that the derivee a gauche of $\Psi_{M/L}$ en $v_L(t)$ is
\begin{equation}\label{eq:3.2.1}
\frac{[G[0]:G[t]]}{[H[0]:H[t]]} = \frac{[G[0]:H[0]]}{[G[t]:H[t]]}.
\end{equation}
As one has $H[t] = G[t] \cap H$, the indice to the denominateur vaut again $[G:G[t]H]$. The derivee of $\Psi_{M/L}$ croit therefore with $t$, and $\Psi_{M/L}$ is convexe, of or (i). the indice \eqref{eq:3.2.1} is constant, of value $e = [G[0]:H[0]]$, of the that $G[t]$ is inclus in $H$. As $G[t]$ is distingue in $G$, that signifie that $G[t]$ agit trivialement in the representation of permutation of $G$ on $G/H$. But this derniere is equivalent a the representation of $G$ on $\Hom_L(M, N)$. One deduces that $H$ contient $G[t]$ if and only if one has $v_L(t) > \alpha(M/L)$ cf.~0.6; that proves (ii) and the premiere assertion of (iii). Enfin, puisque $G[t]$ is a $p$-group for $t$ not zero, the indice $[G[t]:H[t]]$ is a power of $p$, not trivial si $t < \alpha(M/L)$, this which proves the other assertions of (iii).
\end{proof}

\subsection{}
\label{sec:3.3}

Let $L$ a finite extension separable of $K$. One denotes $e = e(L/K)$ the indice of ramification and $f = f(L/K)$ the degree residuel. one has $[L:K] = ef$.

For $y \in L$ of valuation $v_L(y)$ sufficiently grande, one has
\[
N_{L/K}(E(y)) = E(\Tr_{L/K}(y)) \quad \text{et} \quad N_{L/K}(1+y) = 1 + \Tr_{L/K}(y) \pmod{U_K(v_L(y)/e)}.
\]

These formules are verifiees for $v_L(y) > \alpha(L/K)$. We shall preciser this result.

\subsection{}
\label{sec:3.4}

Disons that a element $t$ of $T$ is grand rel.~$L/K$ si the one has $v_L(t) > \Psi_{L/K}(\alpha(L/K))$. Si $t$ is grand, the image of $A_L(v_L(t))$ by $N_{L/K}$ (resp.~$\Tr_{L/K}$) is contenue in $A_K(v_K(t))$ (resp.~$A_K(v_K(t))$). One can therefore, for $t$ grand, considerer the maps
\[
N: A_L(v_L(t))/A_L(v_L(t)+1) \to A_K(v_K(t))/A_K(v_K(t)+1)
\]
and
\[
T: A_L(v_L(t))/A_L(v_L(t)+1) \to A_K(v_K(t))/A_K(v_K(t)+1)
\]
deduites of $N_{L/K}$ and $\Tr_{L/K}$.

In the cas particulier of a extension moderee, $t$ is grand of the that $t > 0$, and $N$ relie the valuation of $x \in \cO_L$ a celle of son image in $\cO_K$.

We we proposons of comparer $N_{L/K}(1+y)$ a $1 + \Tr_{L/K}(y)$ for $y$ in $L$, of grande valuation rel.~$L/K$. By example, for every integer $n > \alpha(L/K)$, one can of after 3.3 considerer the diagram suivant, or the fleches $N$ and $T$ are deduites of $N_{L/K}$ and $\Tr_{L/K}$ respectivement, and the fleches verticales deduites of $y \mapsto E(y)$:
\begin{equation}\label{eq:3.4.1}
\begin{tikzcd}
A_M(v_M(n))/A_M(v_M(n)+1) \arrow[r] \arrow[d] & A_L(v_L(n))/A_L(v_L(n)+1) \arrow[d] \\
U_M(v_M(n))/U_M(v_M(n)+1) \arrow[r] & U_L(v_L(n))/U_L(v_L(n)+1)
\end{tikzcd}
\end{equation}
We verrons that this diagram is commutative. The raisonnement of 3.3 (resp. the proposition 3.5 ci-dessous) the shows lorsque $\Psi_{L/K}(n) > (e/2)n$ (resp. $\Psi_{L/K}(t) > (e/p)n$, this which, by the cas (iii) of the proposition 3.2, implique $n > \alpha(L/K)$).

One obtains of the results more precise en remplacant the function $y \mapsto 1+y$ by the function $E$. S'it the desire, the reader pourra repasser to the prime point of vue en posant $1+y = E(z)$, or $z$ is of valuation to the less $2v(y)$.

The proposition suivante, of proof tres simple, will be precisee by 3.8, proves by devissage and reduction to the cas of a extension cyclique of degree $p$. This amelioration of 3.5 ne we servira that in the proofs of 4.2 (via the commutativity of the diagram 3.4.1) and 4.15.3.

\begin{proposition}\label{prop:3.5}
Let $y$ a element of the ideal maximal of $\cO_L$. one has
\[
N_{L/K}(Ey) = E(\Tr_{L/K}(y)) \pmod{*} U_K((p/e) v_L(y)).
\]
\end{proposition}

\begin{proof}
Plongeant $L$ in a extension galoisienne $M$ of $K$, one can ecrire the norme (resp.~the trace) of a element $x$ of $L$ as product (resp.~sum) of conjugues $x^\sigma$ of this element. One trouve then that $N_{L/K}(Ey) = \prod E(y^\sigma)$ is congru a $E(\Tr_{L/K}(y))$ modulo of the elements of $M$ of valuation to the less celle of $y^p$. The proposition en resulte aussitot.
\end{proof}

\begin{proposition}\label{prop:3.6}
Let $\chi$ a quasi-character sauvagement ramifie of $K^*$, verifiant $\Psi_{L/K}(\alpha(\chi)) > (e/p) \alpha(\chi)$. Let $a$ a element of $K$ such that the one ait $\chi(Ey) = \psi(ay)$ for $y$ in $K$ of valuation more grande that $\alpha(\chi)/p$, and $a'$ a element of $L$ such that the one ait $\chi \circ N_{L/K}(Ey) = \psi \circ \Tr_{L/K}(a'y)$ for $y$ in $L$ of valuation more grande that $\alpha(\chi \circ N_{L/K})/p$. Then one has
\[
a = a' \pmod{*} U_L(\Psi_{L/K}(\alpha(\chi)) - (e/p) \alpha(\chi)).
\]
\end{proposition}

Notons of abord that the hypothese implique $\alpha(\chi) > \alpha(L/K)$ and $\alpha(\chi \circ N_{L/K}) = \Psi_{L/K}(\alpha(\chi))$. Prenons a element $y$ of $L$, of valuation $v_L(y) > e\alpha(\chi)/p$. The proposition 3.5 given then
\[
\chi \circ N_{L/K}(Ey) = \chi(E(\Tr_{L/K}(y))).
\]

But, of after 3.3, one has $\alpha(\chi \circ N_{L/K}) = \Psi_{L/K}(\alpha(\chi)) \leq e\alpha(\chi)$. For $v_L(y) > \alpha(\chi \circ N_{L/K})/p$ and $\chi \circ N_{L/K}(Ey) = \psi \circ \Tr_{L/K}(a y)$. Of 3.3 one deduces also that the one has $\Psi_{L/K}(v_L(\Tr_{L/K}(y))) \geq v_L(y)$, that is
\[
v_K(\Tr_{L/K}(y)) + \Psi_{L/K}(\alpha(\chi)) - e\alpha(\chi) \geq v_L(y)
\]
or
\[
v_K(\Tr_{L/K}(y)) \geq \alpha(\chi)/p + (e\alpha(\chi) - \Psi_{L/K}(\alpha(\chi))),
\]
and
\[
v_K(\Tr_{L/K}(y)) > \alpha(\chi)/p.
\]
One en tire the equality $\chi(E(\Tr_{L/K}(y))) = \psi(a \Tr_{L/K}(y))$. For every element $y$ of $K$ verifiant $v_L(y) > e\alpha(\chi)/p$, one has therefore
\[
\psi \circ \Tr_{L/K}((a-a')y) = 1,
\]
this which implique
\[
v_L(a'-a) \geq -n(\psi \circ \Tr_{L/K}) - (e\alpha(\chi)/p) - 1.
\]
As $a'$ is of valuation $-n(\psi \circ \Tr_{L/K}) - \Psi_{L/K}(\alpha(\chi)) - 1$, one has bien
\[
a = a' \pmod{*} U_L(\Psi_{L/K}(\alpha(\chi)) - (e/p) \alpha(\chi)). \qedhere
\]

\begin{corollary}\label{cor:3.7}
Let $\chi$ a quasi-character sauvagement ramifie of $K^*$ and $\eta$ a quasi-character of $L^*$. One assumed that, posant $u = \sup(\alpha(\chi), \Psi_{L/K}(\alpha(L/K)))$, one ait $\Psi_{L/K}(u) < (1-1/p) \alpha(\chi)$. One defines $a$ and $a'$ as en \ref{prop:3.6}, and one has then: $v(a) = v(a')$.
\end{corollary}

the hypothese faite on $u$ we given of abord
\[
\alpha(\chi) > \Psi_{L/K}(u) \geq \alpha(L/K),
\]
\[
\alpha(\chi) + (u - \Psi_{L/K}(\alpha(\chi)))/e < (1 - (1/p)) \alpha(\chi),
\]
then
\[
u < \Psi_{L/K}(\alpha(\chi)) - e\alpha(\chi)/p,
\]
of or $u < \Psi_{L/K}(\alpha(\chi)) - e\alpha(\chi)/p$, and one applique \ref{prop:3.6}.

\subsection{}
\label{sec:3.8}

The graphique suivant visualise the proposition ci-dessous, and sa relation with \ref{prop:3.5}. In the graphique, $T_L$ and $T_K$ are representes a of the echelles differentes (unites $e$ and $1$).

\begin{proposition}\label{prop:3.8}
Let $t$ a element of $T$ grand rel.~$L/K$. Let $y$ a element of $A_L(v_L(t))$. one has $\Tr_{L/K}(y) \in A_K(v_K(t))$ and
\[
N_{L/K}(Ey) = E(\Tr_{L/K}(y)) \pmod{*} U_K(\alpha(L/K) + p(v_K(t) - \alpha(L/K))).
\]
\end{proposition}

Notons $R_{L/K}$ the function lineaire of pente $e/p$, defined for $t = \alpha(L/K)$ and co\"incidant with $\Psi_{L/K}$ en $\alpha(L/K)$; notons $R_{L/K}^{-1}$ the function inverse. the assertion of \ref{prop:3.8} signifie that si the valuation of $y$ is grande rel.~$L/K$, one has
\[
N_{L/K}(Ey) = E(\Tr_{L/K}(y)) \pmod{*} U_K(R_{L/K}^{-1}(v_L(y))).
\]

\begin{proof}[Proof]
Supposons of abord that the extension $L/K$ let cyclique of degree $p$, totalement ramifiee. En this cas, etudie en [2], V \S 3, $\Psi_{L/K}$ is of pente $1$ for $t < \alpha(L/K)$, of pente $p$ for $t > \alpha(L/K)$, and the one has $R_{L/K}^{-1}(t) = t$. As en \ref{prop:3.5}, one trouve $N_{L/K}(Ey) \equiv E(\Tr_{L/K}(y)) \pmod{*} U_K(pv_L(y)) \cap K^*$. As $U_K(pv_L(y)) \cap K^*$ n'is other that $U_K(v_L(y))$, the assertion of \ref{prop:3.8} is therefore vraie.

Supposons the proposition \ref{prop:3.8} vraie for deux extensions $M/L$ and $L/K$, and prouvons-the for the extension $M/K$. Puisque the one has $\Psi_{M/K} = \Psi_{M/L} \circ \Psi_{L/K}$ and that the functions $\Psi$ are convexes, $v_L(t)$ is a point of discontinuite of $\Psi_{M/K}$ if and only if c'en is a for $\Psi_{M/L}$, or that $v_L(t)$ en is a for $\Psi_{L/K}$. In particular, $t$ is grand rel.~$M/K$ if and only if it the is rel.~$M/L$ and rel.~$L/K$. En this cas, si $y$ is a element of $A_M(v_M(t))$, one has
\[
N_{M/K}(Ey) = N_{L/K}(N_{M/L}(Ey)) = N_{L/K}(E(\Tr_{M/L}(y))) \equiv E(\Tr_{L/K}(\Tr_{M/L}(y))) = E(\Tr_{M/K}(y)),
\]
the congruences being prises $\mod^*\, U_K(R_{L/K}^{-1}(v_L(t)))$. It s'agit therefore of prouver the inegalites
\[
R_{M/L}^{-1} \circ R_{L/K}^{-1} \geq R_{M/K}^{-1} \quad \text{et} \quad R_{M/L}^{-1} \circ R_{L/K}^{-1} \geq R_{M/K}^{-1}
\]
en $v_M(t)$, for $t$ grand rel.~$M/K$. In this domaine, these functions are lineaires similarly pente, and it is enough to verify the inegalites en $v_M(t)$, quand $v_M(t)$ vaut $\alpha(M/K)$. En this point, one has $R_{M/K}^{-1} = \Psi_{M/K} = \Psi_{M/L} \circ \Psi_{L/K}$ and $R_{M/K}^{-1} \leq R_{M/L}^{-1} \circ R_{L/K}^{-1}$ ainsi that $R_{M/K}^{-1} \geq R_{M/L}^{-1} \circ R_{L/K}^{-1}$, of or the assertion of \ref{prop:3.8} for $M/K$.
\end{proof}

\subsection{}
\label{sec:3.11}

Grimpant the long of a tour, one deduces of this which precede that the proposition \ref{prop:3.8} is vraie for every extension $L/K$ which se plonge in a extension galoisienne $\tilde{L}/K$, totalement ramifiee and of Galois group a $p$-group. Indeed, for every under-group $H$ of a $p$-group $G$, there exists a sequence of under-groups $H = H_0 \triangleleft H_1 \triangleleft \cdots \triangleleft H_n = G$, or $H_i$ is distingue in $H_{i+1}$ and or the quotients $H_{i+1}/H_i$ are of ordre $p$.

Passons to the cas general. Let $t_0$ and $t$ deux elements of $T$ such that the one ait $v_K(t_0) = \alpha(L/K)$ and $t > t_0$. Prenons $y$ in $L^*$, of valuation $v_L(t)$. It s'agit of prouver the congruence
\[
N_{L/K}(Ey) = E(\Tr_{L/K}(y)) \pmod{*} U_K(v_K(t_0) + p(v_K(t) - v_K(t_0))).
\]

Let $K'$ a extension moderement ramifiee of $K$. Decomposons $K' \otimes_K L$ en product of the field $L_i$. Notant $y_i$ the image in $L_i$ of the element $y$ of $L$, one has
\[
N_{L/K}(Ey) = \prod_i N_{L_i/K'}(Ey_i).
\]
The extensions $L_i/L$ are moderees. The functions $\Psi_{L_i/L}$ and $\Psi_{L_i/K'}$ are therefore lineaires. Puisque $\Psi_{L_i/K}$ n'is other that $\Psi_{L_i/K'} \circ \Psi_{K'/K}$ and that $\Psi_{K'/K}$ is lineaire, one has $v_{K'}(t_0) = \alpha(L_i/K')$. Supposons the proposition vraie for the extensions $L_i/K'$. Each $y_i$ being of valuation $v_{L_i}(t)$, one has then
\[
N_{L_i/K'}(Ey_i) = E(\Tr_{L_i/K'}(y_i)) \pmod{*} U_{K'}(v_{K'}(t_0) + p(v_{K'}(t) - v_{K'}(t_0))),
\]
and
\[
\prod_i E(\Tr_{L_i/K'}(y_i)) \equiv E(\sum_i \Tr_{L_i/K'}(y_i)) = E(\Tr_{L/K}(y)) \pmod{*} U_K(p(v_K(t) - v_K(t_0))).
\]
The deux membres of this derniere congruence being in $K$, one has
\[
N_{L/K}(Ey) = E(\Tr_{L/K}(y)) \pmod{*} U_K(v_K(t_0) + p(v_K(t) - v_K(t_0))).
\]
But, as one sait, one can choisir $K'$ in such a way that the extensions $L_i/K'$ satisfy the condition of 3.11. The proposition \ref{prop:3.8} is vraie for the $L_i/K'$ therefore for $L/K$.

\begin{corollary}\label{cor:3.12}
\begin{enumerate}[label=(\roman*)]
\item For every integer $n > \alpha(L/K)$, the diagram \eqref{eq:3.4.1} is commutative.
\item Let $\chi$ a quasi-character sauvagement ramifie of $K^*$, verifiant $\alpha(\chi) > \alpha(L/K)$. Definissons $t_0 \in T$ by the equality $v_K(t_0) = \alpha(L/K)$ and $t \in T$ by $v_K(t) = \alpha(\chi)$. Let $a$ a element of $K$ such that the one ait $\chi(Ey) = \psi(ay)$ for $y \in K$ verifiant $pv(y) > \alpha(\chi)/p$ and $a'$ a element of $L$ such that the one ait $\chi \circ N_{L/K}(Ey) = \psi \circ \Tr_{L/K}(a'y)$ for $y \in L$ verifiant $pv_L(y) > \alpha(\chi \circ N_{L/K})/p$. These conditions determinent $a$ $\mod^*\, U_K((1-1/p)(\alpha(\chi) - v_K(t_0)))$, $a'$ $\mod^*\, U_L((1-1/p)(\alpha(\chi) - v_L(t_0)))$, and the one has
\[
a = a' \pmod{*} U_L((1-1/p)(v_L(t) - v_L(t_0))).
\]
\end{enumerate}
\end{corollary}

The part (i) is immediate, si one remark that the one has $\alpha(L/K) + p(n - \alpha(L/K)) > en$. For prouver (ii) prenons a element $y$ of $L$ verifiant $v_L(y) > v_L(t) - (e/p)(v_L(t) - v_L(t_0))$; one has therefore $R_{L/K}^{-1}(v_L(y)) > v_K(t) = \alpha(\chi)$ and, therefore $\chi \circ N_{L/K}(Ey) = \chi(E(\Tr_{L/K}(y)))$. one has also
\[
v_K(\Tr_{L/K}(y)) \geq v_K(t_0) + (1/p)(v_L(y) - v_L(t_0)) > v_K(t_0) + (1/p)(v_L(t) - v_L(t_0)).
\]
One deduces the equality $\chi \circ N_{L/K}(Ey) = \psi \circ \Tr_{L/K}(ay)$, that is that $a$ satisfies the property caracteristique of $a'$.


\section{The cas general}
\label{sec:4}

\subsection{}
\label{sec:4.1}

Let $V$ a representation virtuelle of $W$, and let $\gamma$ a element of $K^*$ of valuation $a(V) + \dim V \cdot n(\psi)$.
Si $\eta$ is a quasi-character not ramifie of $K^*$, one has
\begin{equation}\label{eq:4.1.1}
\varepsilon(\eta \otimes V, \psi, dx) = \varepsilon(V, \psi, dx) \cdot \eta(\gamma).
\end{equation}
More generally, si $W$ is a representation virtuelle of $W$, not ramifiee and of dimension 0, one has (cf.~[1] 5.5.3)
\begin{equation}\label{eq:4.1.2}
\varepsilon(W \otimes V, \psi) = \det W(\gamma).
\end{equation}

The results annonces in the introduction on the calcul of $\varepsilon(W \otimes V, \psi)$, quand $W$ is of dimension 0 and less ramifie that $V$, generalisent this formule \eqref{eq:4.1.2} and, as the reader the verifiera aisement, s'y reduisent si the space $V^P$ of the points fixes by $P$ is not trivial, this which, with the notations of the introduction, impose $t=0$. For prouver these generalisations, we supposerons therefore desormais that the space $V^P$ is trivial, i.e.~that the representation $V$ n'a not of constituant modere.

The scheme of proof of these results is the suivant. One assumed of abord $V$ of dimension 1. Si $W$ is of the form $\eta - 1$, or $\eta$ is a quasi-character of $K^*$, the result cherche resulte of the chapitre 1. For passer of the to the cas general, one ecrit, grace to the chapitre 2, $V$ as combinaison lineaire of induites of quasi-characters $\chi_i$ of extensions $L_i$ of $K$, and one se ramene, by induction, to the cas of dimension 1. The results for $W$ of dimension 0, trivial on $W(t)$ but not on $W(t/2)$, s'en deduisent by passage a the limite.

\begin{theorem}\label{thm:4.2}
Let $V$ a representation virtuelle of $W$, without constituant modere. There exists a element $\gamma$ of $K^*$, determine in such a way unique $\mod^*\, U(1)$, such that for every representation virtuelle $W$ of $W$, of dimension 0 and verifiant $\beta(W) < \alpha(V)$, one ait
\begin{equation}\label{eq:4.2.1}
\varepsilon(W \otimes V, \psi) \equiv \det W(\gamma) \pmod{*} \bQ_p/\bZ_p(1).
\end{equation}
The valuation of $\gamma$ is $a(V) + (\dim V) \cdot n(\psi)$.
\end{theorem}

\begin{proof}[Proof of the unicite]
the unicite of $\gamma$ is clear: prenant $W = \eta - 1$, or $\eta$ is a quasi-character of $K^*$, one sees that $\gamma$ is determine $\mod^*\, U(1)$ by the values of $\varepsilon((\eta-1) \otimes V, \psi)$ for $\eta$ modere. Indeed, the formule \eqref{eq:4.2.1} determine $\eta(\gamma)$ for every character $\eta$ of $U(0)/U(1)$, puisque this group is of ordre prime a $p$.
\end{proof}

\begin{proof}[Proof: cas $V$ of dimension 1]
Gardons the hypotheses and notations of \ref{thm:4.2}, and supposons en outre $V$ of dimension 1, corresponding to the quasi-character (sauvagement ramifie) $\chi$ of $K^*$. Let $a$ the element of $K^*$, uniquement defined $\mod^*\, U(1)$, such that the one ait $\chi(1+y) = \psi(ay)$ for $y \in K$ verifiant $v_K(y) \geq \alpha(\chi)$. Then $\gamma = a^{-1}$ satisfies the formule \eqref{eq:4.2.1}.

By hypothese, the constituants of $W$ are trivial on $W(\alpha(\chi))$. of after the variante ([1], 1.5) of the theorem of Brauer en dimension 0, $W$ s'ecrit as combinaison lineaire of representations of the form $\Ind_H^W(\cX - 1)$, or $H$ is a under-group of $W$ contenant $W(\alpha(\chi))$, and $\cX$ a quasi-character of $H$ trivial on $W(\alpha(\chi))$. The deux membres of \eqref{eq:4.2.1} are multiplicatifs en $W$, one can therefore supposer $W$ lui-same of this form $W = \Ind_H^W(\cX - 1)$. Let $L$ the finite extension of $K$ fixee by $H$ and $\eta$ the quasi-character of $L^*$ corresponding a $\cX$.

the hypothese that $\cX$ let trivial on $W(\alpha(\chi))$ se traduit by the inegalites
\[
\alpha(L/K) < \alpha(\chi) \quad \text{et} \quad \alpha(\eta) < \Psi_{L/K}(\alpha(\chi)).
\]
\end{proof}

\subsection{}
\label{sec:4.4}

one has $W \otimes V = \Ind_H^W(\Ind_H^H((\cX-1) \otimes V))$ and $V|_H$ s'identifie to the quasi-character $\chi \circ N_{L/K}$ of $L^*$. The compatibility of the factors $\varepsilon$ a the induction, en dimension 0, we fournit then the equality
\[
\varepsilon(W \otimes V, \psi) = \varepsilon((\eta-1) \cdot \chi \circ N_{L/K}, \psi \circ \Tr_{L/K}).
\]
the hypothese $\alpha(L/K) < \alpha(\chi)$ assure that the one has $\alpha(\chi \circ N_{L/K}) = \Psi_{L/K}(\alpha(\chi))$ and that, for $y \in L$ of valuation $\Psi_{L/K}(\alpha(\chi))$, one has (cf.~\ref{cor:3.12}, (i))
\[
N_{L/K}(Ey) = E(\Tr_{L/K}(y)) \pmod{*} U(\alpha(\chi)+1).
\]
Therefore
\[
\chi \circ N_{L/K}(Ey) = \chi(E(\Tr_{L/K}(y))) = \psi(a \Tr_{L/K}(y)) = \psi \circ \Tr_{L/K}(ay).
\]
the hypothese $\alpha(\eta) < \Psi_{L/K}(\alpha(\chi)) = \alpha(\chi \circ N_{L/K})$ permet of appliquer the corollary \ref{cor:1.3}(i); one deduces the equality
\[
\varepsilon((\eta-1) \cdot \chi \circ N_{L/K}, \psi \circ \Tr_{L/K}) \equiv \psi(a^{-1}) \pmod{*} \bQ_p/\bZ_p(1).
\]
For every representation virtuelle $R$ of $W(K/L)$, of dimension 0, one has
\[
\det \Ind_H^W(R) = \det R|_{H^{\text{ab}}} \circ \text{transfert}
\]
(cf.~[2], XI \S 3 and [1] 1.2, which reproduit a result of P.X.~Gallagher, \textit{Determinant of representations of finite groups}, Abh.~Math.~Sem.~A.~Hamburg, 28 (1965)). Ici, one trouve that the one has $\det W = \eta|_{K^*}$ and the equality
\[
\varepsilon(W \otimes V, \psi) \equiv \eta(a^{-1}) = \det W(a^{-1}) \pmod{*} \bQ_p/\bZ_p(1)
\]
proves the theorem for $V$ of dimension 1.

\subsection{}
\label{sec:4.5}

Prouvons maintenant \ref{thm:4.2} in the cas general. One can supposer $V$ irreducible.

Let $G$ the quotient of $W$ by the under-group of $I$ which agit trivialement on $V$. The group $G(\alpha(V))$ is the dernier group of ramification not zero of $G$; it is therefore abelien, and distingue in $G$. It agit on $V$ without vecteur fixe not zero (otherwise, the space of the vecteurs fixes, stable under $W$ serait not trivial, and the representation $V$ serait reductible). Appliquant the variante \ref{prop:2.3}, one sees that $V$ is a combinaison lineaire $\sum n_i \Ind_{H_i}^W(\chi_i)$ or the under-group $H_i$ of $G$ contient $G(\alpha(V))$ and or the character $\chi_i$ of $H_i$ is not trivial on $G(\alpha(V))$. A $H_i$ correspond a extension $L_i$ of $K$, and a $\chi_i$ a quasi-character of $L_i^*$, again note $\chi_i$. Let $W_i$ the restriction of $W$ a $W(K/L_i)$. one has $\beta(W_i) \leq \Psi_{L_i/K}(\beta(W))$ and $\alpha(\chi_i) = \Psi_{L_i/K}(\alpha(V))$, therefore also $\beta(W_i) < \alpha(\chi_i)$. Si $a_i$ is a element of $L_i$ such that the one ait $\chi_i(1+y) = \psi \circ \Tr_{L_i/K}(a_i y)$ for $y \in L_i$ verifiant $v_{L_i}(y) \geq \alpha(\chi_i)$, one has
\begin{align*}
\varepsilon(W \otimes V, \psi) &= \prod_i \varepsilon(W_i \otimes V_i, \psi \circ \Tr_{L_i/K})^{n_i} \\
&= \prod_i \det W_i(a_i^{-1})^{n_i} \quad \text{par le cas de dim.~1} \\
&= \prod_i \det W(N_{L_i/K}(a_i^{-1}))^{n_i}.
\end{align*}
the element $\gamma = \prod N_{L_i/K}(a_i^{-1})^{n_i}$ satisfies \eqref{eq:4.2.1}. C.Q.F.D.

\begin{theorem}\label{thm:4.6}
Let $V$ a representation virtuelle of $W$, without constituant modere. There exists a element $\gamma$ of $K^*$, defined in such a way unique $\mod^*\, U(\alpha(V)/2)$, such that for every representation virtuelle $W$ of $W$, of dimension 0 and verifiant $\beta(W) < \alpha(V)/2$, one ait
\begin{equation}\label{eq:4.6.1}
\varepsilon(W \otimes V, \psi) = \det W(\gamma).
\end{equation}
This element $\gamma$ is congru a celui of \eqref{eq:4.2.1} $\mod^*\, U(1)$.
\end{theorem}

The proof of this theorem is parallele a the precedente. the unicite of $\gamma$ is clear, it vaut deja lorsqu'one se restreint a prendre $W$ of the form $\eta - 1$ or $\eta$ is a quasi-character quelconque of $K^*$, trivial on $U(\alpha(V)/2)$. The compatibility with \eqref{eq:4.2.1} resulte of \ref{thm:4.2}(iii).

\begin{proof}[Proof: cas $V$ of dimension 1]
Gardons the hypotheses and notations of the theorem \ref{thm:4.6}, and supposons en outre $V$ of dimension 1, corresponding to the quasi-character $\chi$ of $K^*$. Let $a$ a element of $K^*$, defined in such a way unique $\mod^*\, U(\alpha(\chi)/2)$, such that the one ait $\chi(Ey) = \psi(ay)$ for $y \in K$ verifiant $v(y) > \alpha(\chi)/2$. Then $\gamma = a^{-1}$ satisfies \eqref{eq:4.6.1}.

Si the one has $\beta(W) < \alpha(\chi)/2$, one has also $a(\det W) < \alpha(\chi)/2$, and the validite of \eqref{eq:4.6.1} ne depend not of the choix of $a$. One simplifiera the references en prenant $a$ such that the one ait $\chi(Ey) = \psi(ay)$ for $y \in K$ verifiant $v(y) > \alpha(\chi)/p$.

Of the same maniere that en 4.4, one se ramene a supposer $W$ of the form $\Ind_H^W(\cX-1)$, or the under-group $H$ of $W$ correspond a extension $L$ of $K$, and the quasi-character $\cX$ of $H$ a quasi-character $\eta$ of $L^*$, verifiant:
\[
\alpha(L/K) < \alpha(\chi)/2 \quad \text{et} \quad \alpha(\eta) < \Psi_{L/K}(\alpha(\chi)/2).
\]
The premiere inegalite assure that the one has $\alpha(\eta \circ N_{L/K}) = \Psi_{L/K}(\alpha(\chi))$ and the convexite of $\Psi_{L/K}$ we given $\Psi_{L/K}(\alpha(\chi)/2) \leq \Psi_{L/K}(\alpha(\chi))/2$. One can therefore appliquer the corollary \ref{cor:1.3}(ii) a $\eta \circ N_{L/K}$; si $a'$ is a element of $L^*$ such that the one ait
\[
\chi \circ N_{L/K}(Ey) = \psi \circ \Tr_{L/K}(a' y)
\]
for $y \in L^*$ verifiant $v(y) > \alpha(\chi \circ N_{L/K})/2$, one has
\[
\varepsilon(W \otimes V, \psi) = \varepsilon((\eta-1) \cdot \chi \circ N_{L/K}, \psi \circ \Tr_{L/K}) = \eta(a'^{-1}).
\]
The corollary \ref{cor:3.7} we assure the equality $v(a) = v(a')$, and one conclut as en 4.4.
\end{proof}

\subsection{}
\label{sec:4.8}

One passe of the lemma to the theorem as en 4.5. One se ramene a supposer $V$ irreducible, one ecrit $V$ under the form $\sum n_i \Ind_{H_i}^W(\chi_i)$, one defines a element $a_i$ of $L_i$ en exigeant the equality $\chi_i(1+y) = \psi \circ \Tr_{L_i/K}(a_i y)$ for $v(y) > \alpha(\chi_i)/2$, and one trouve that the element $\gamma = \prod N_{L_i/K}(a_i^{-1})^{n_i}$ satisfies \eqref{eq:4.6.1}. The convexite of $\Psi_{L_i/K}$ garantit the inegalite requise:
\[
\beta(W_i) \leq \Psi_{L_i/K}(\beta(W)) < \Psi_{L_i/K}(\alpha(V)/2) \leq \alpha(\chi_i)/2.
\]

\begin{theorem}\label{thm:4.9}
Let $V$ a representation virtuelle of $W$, without constituant modere. Let $(W_i)_{i \in I}$ a family finite of representations virtuelles of dimension 0 of $W$. One assumed that the one has $\beta(W_i) < \alpha(V)$ for every indice $i$ and $\sum_{i \in I}(\alpha(V) - \beta(W_i)) > \alpha(V)$. Then one has
\[
\varepsilon\biggl(\bigl(\bigotimes_i W_i\bigr) \otimes V\biggr) = 1.
\]
\end{theorem}

\begin{remark}
the hypothese assure that there is to the less deux representations $W_i$. The determinant of $(\bigotimes_i W_i) \otimes V$ is therefore trivial, this which justifie the omission of $\psi$ under the signe $\varepsilon$.
\end{remark}

\begin{proof}
Si $W_i$ is of the form $\Ind_{H_i}^W(\eta_i - 1)$ or $L_i$ is a finite extension of $K$ and $\eta_i$ a quasi-character of $L_i^*$, verifiant
\[
\alpha(\eta_i) < \Psi_{L_i/K}(\beta(W_i)) \quad \text{et} \quad \alpha(L_i/K) \leq \beta(W_i),
\]
one has
\begin{equation}\label{eq:4.9.1}
\varepsilon\biggl(\bigl(\bigotimes_i W_i\bigr) \otimes V\biggr) = \varepsilon\biggl(\bigl(\bigotimes_i (W_i|_{W(K/L_i)}) \otimes (V|_{W(K/L_i)})\bigr), \psi \circ \Tr_{L_i/K}\biggr),
\end{equation}
and moreover
\[
\beta(W_i|_{W(K/L_i)}) \leq \Psi_{L_i/K}(\beta(W_i))
\]
\[
\alpha(V|_{W(K/L_i)}) = \Psi_{L_i/K}(\alpha(V)).
\]
The convexite of $\Psi_{L_i/K}$ assure therefore that the hypothese of the theorem is again satisfaite for the representation (of $W(K/L_i)$) apparaissant in the membre of line of \eqref{eq:4.9.1}.

Similarly, si $V$ is of the form $\Ind_{H_i}^W(\chi_i)$, or the quasi-character $\chi_i$ of $L_i^*$ satisfies $\alpha(\chi_i) = \Psi_{L_i/K}(\alpha(V))$, one has
\begin{equation}\label{eq:4.9.2}
\varepsilon\biggl(\bigl(\bigotimes_i W_i\bigr) \otimes V\biggr) = \varepsilon\biggl(\bigotimes_i (W_i|_{W(K/L_i)}) \otimes \chi_i, \psi \circ \Tr_{L_i/K}\biggr),
\end{equation}
and the hypothese of the theorem is again satisfaite for the representation apparaissant to the membre of line of \eqref{eq:4.9.2}.

The variantes of the theorem of Brauer ramenent therefore the theorem a the proposition \ref{prop:1.7}.
\end{proof}

\subsection{}
\label{sec:4.10}

\textbf{Question.} Let $V$ a representation of $W$, without constituant modere. Let $(G_i)_{i \in I}$ a family finite of quotients of $W$ by of the under-groups ouverts of $I$ and, for each indice $i$, $f_i$ the dernier saut of the filtration of $G_i$ by the groups of ramification, en numerotation superieure. One assumed that for each $i$, one has $f_i < \alpha(V)$. Let $(n_i)_{i \in I}$ a family of integers positifs or nuls, and for each indice $i$, a representation virtuelle $W_i$ of $G_i$, which se trouve in the $n_i$-under-space of the $\gamma$-filtration of the ring of the representations of $G_i$. For $n_i = 1$ (resp.~$n_i = 2$) that signifie that $W_i$ is of dimension 0 (resp.~of dimension 0 and of determinant trivial).

Si $\sum_{i \in I} n_i(\alpha(V) - \beta(W_i)) > \alpha(V)$, can-one conclure that $\varepsilon((\bigotimes_i W_i) \otimes V)$ vaut 1?

\begin{remark}
Si all the $n_i$ valent 1, c'is the theorem \ref{thm:4.9}. Si $|I|$ vaut 1 and si $n_i$ vaut 2, c'is the theorem \ref{thm:4.6}.
\end{remark}

\begin{lemma}\label{lem:4.11}
Every representation irreducible of the group of inertie sauvage of which the class of isomorphie is fixe under $W$ se prolonge en a representation of $W$.
\end{lemma}

The group of inertie $I$ is a extension
\begin{equation}\label{eq:4.11.1}
1 \to P \to I \to \bZ_\ell(1) \to 1.
\end{equation}
Let $\sigma_0$ a generateur of $\bZ_\ell(1)$, $\tilde{\sigma}_0$ a lifting of $\sigma_0$ in $I$, $1(\neq p)$ the element of $\hat{\bZ} = \prod_\ell \bZ_\ell$ of composante in $\bZ_\ell$ egale a 1 (resp.~a 0) for $\ell \neq p$ (resp.~$\ell = p$), and posons $c = \tilde{\sigma}_0^{1(\neq p)}$. This element defines a scindage of the extension \eqref{eq:4.11.1}.

The group $W$ is extension of $\bZ$ by $I$. The choix of a lifting $F$ of $1 \in \bZ$ scinde this extension; celui of $F$ and $\sigma_0$ identifie $W$ to the product semi-direct itere
\[
W \simeq \bZ \ltimes (\bZ_\ell \ltimes P).
\]
For prolonger a $W$ a representation $\rho: P \to GL(V)$, it is enough therefore of se donner the deux automorphismes $R = \rho(\sigma_0)$ and $S = \rho(F)$ of $V$ assujettis to the conditions suivantes:
\begin{enumerate}[label=(\alph*)]
\item For $g \in P$, $\rho(\sigma_0 g \sigma_0^{-1}) = R \rho(g) R^{-1}$.
\item $\rho(c)$ is of ordre finite prime a $p$.
\item Notant $\rho_0$ the representation of $I$ which prolonge $\rho$ and for which $R = \rho(\sigma_0)$, one has $\rho_0(F g F^{-1}) = S \rho_0(g) S^{-1}$ for $g \in I$.
\end{enumerate}

Supposons $V$ irreducible, of class of isomorphie fixe by $I$. There exists then $R_0$ verifiant (a). Let $n$ a integer $\geq 1$, prime a $p$, such that $\sigma_0^n$ centralise the quotient finite $\rho(P)$ of $P$. The representation $\rho$ being irreducible, the automorphism $R_0^n$ of $(V, \rho)$ is a scalaire $\lambda$. Let $\lambda^{1/n}$ a racine $n$\textsuperscript{eme} of $\lambda$ in $\bC^*$; posons $R = \mu R_0$, or $\mu^{n-1} = \lambda^{1/n}$. Then $R$ satisfies (a) and $R^n = 1$. One can choisir $S$ verifiant (c), and remplacant to the need $S$ by a power $S^{1 + Mp^k}$ with $M$ suitable, one obtains (b). C.Q.F.D.


\subsection{}
\label{sec:4.13}

Let $V$ a representation virtuelle of $W$. Decomposons the restriction of $V$ a $P$ en sum of composantes isotypiques:
\[
V|_P = \sum_{\nu \in \hat{P}/W} V(\nu),
\]
or $V(\nu)$ is the composante isotypique corresponding a class of isomorphie $\nu$ of representations irreductibles of $P$, and the sum porte on the orbites of $W$ in $\hat{P}$. One can supposer that $\nu$ apparait in a constituant of $V$. Si $\nu$ is the representation trivial, one sets $V(\nu) = V^P$.

Let $W(\nu)$ the fixateur of $\nu$ in $W$. of after the lemma \ref{lem:4.11}, there exists a representation $V_\nu$ of $W(\nu)$ such that $V(\nu) = \Ind_{W(\nu)}^W(V_\nu)$. Posons
\[
V' = \sum_{\nu \neq 1} \Ind_{W(\nu)}^W\bigl((\dim \Hom_P(\nu, V)) \cdot V_\nu\bigr)
\]
and
\[
V'' = \sum_{\nu \neq 1} \Ind_{W(\nu)}^W\bigl((\Hom_P(\nu, V) - (\dim \Hom_P(\nu, V)) \cdot 1) \otimes V_\nu\bigr).
\]
one has $V = V' + V''$. For every representation virtuelle $W$ of $W$, one has
\[
\varepsilon(W \otimes V, \psi) = \varepsilon(W \otimes V', \psi) \cdot \varepsilon(W \otimes V'', \psi).
\]

\begin{theorem}\label{thm:4.13}
With the notations precedentes:
\begin{enumerate}[label=(\roman*)]
\item one has $\varepsilon(W \otimes V'', \psi) = \det W(\gamma'')$, or $\gamma''$ is the element of the theorem \ref{thm:4.6} relative a $V''$.
\item Si $W$ is trivial on $W(t)$, one has $\varepsilon(W \otimes V', \psi) = 1$.
\end{enumerate}
\end{theorem}

\begin{proof}
Prouvons (ii). By definition, $V'$ is induite a partir of representations $V'_\nu$ of $W(\nu)$ of which the restriction a $P$ is scalaire (of dimension $\dim V(\nu)$). one has therefore
\[
\varepsilon(W \otimes V', \psi) = \prod_\nu \varepsilon(W|_{W(\nu)} \otimes V'_\nu, \psi).
\]
This justifie the omission of $\psi$ under the signe $\varepsilon$ in (i). Moreover,
\[
\varepsilon(W \otimes V') = \prod_\nu \varepsilon(W \otimes \Hom_P(V(\nu), V) \otimes V(\nu)).
\]
It is enough to etendre the product to the $\nu$ which apparaissent in a constituant of $V$, each factor is justifiable of \ref{thm:4.9}, and this proves (i).

Assume that $W$ ait a restriction a $P$ zero. Decomposons $W$ en the sum $W = W_1 + W_2$, with
\[
W_1 = \sum_\nu \Ind_{W(\nu)}^W\bigl(\Hom_P(V(\nu), W) \otimes (V(\nu) - \dim V(\nu) \cdot 1)\bigr)
\]
and
\[
W_2 = \sum_\nu \Ind_{W(\nu)}^W\bigl(\Hom_P(V(\nu), W) \otimes (\dim V(\nu) \cdot 1)\bigr).
\]
The representations virtuelles $\Hom_P(V(\nu), W) \otimes V(\nu)$ and $\Hom_P(V(\nu), W) \otimes (\dim V(\nu) \cdot 1)$ of $W(\nu)$ ayant same determinant, one has $\det W = \det W_1$ (cf.~4.4). one has
\[
\varepsilon(W \otimes V, \psi) = \varepsilon(W_1 \otimes V, \psi) \cdot \varepsilon(W_2 \otimes V, \psi).
\]
The representation $W_2 \otimes V$ is of determinant trivial, and the factor $\varepsilon(W_2 \otimes V, \psi)$ se recrit
\[
\varepsilon(W_2 \otimes V) = \prod_\nu \varepsilon\bigl(\Hom_P(V(\nu), W) \otimes (V(\nu) - \dim V(\nu) \cdot 1) \otimes V\bigr).
\]
It is enough to etendre the product to the $\nu$ which apparaissent in the constituants of $W$; each factor is justiciable of \ref{thm:4.9}, and $\varepsilon(W_2 \otimes V) = 1$. Appliquant a $W_1$, one trouve enfin that
\[
\varepsilon(W \otimes V, \psi) = \varepsilon(W_1 \otimes V, \psi) = \det W_1(\gamma) = \det W(\gamma). \qedhere
\]
\end{proof}

\begin{remark}\label{rem:4.14}
For every representation virtuelle $V$ of $W$, posons $V = V' + V''$, with
\[
V' = \sum_\nu \Ind_{W(\nu)}^W\bigl((\dim \Hom_P(\nu, V)) \cdot V_\nu\bigr)
\]
and
\[
V'' = \sum_\nu \Ind_{W(\nu)}^W\bigl((\Hom_P(\nu, V) - (\dim \Hom_P(\nu, V)) \cdot 1) \otimes V_\nu\bigr).
\]
The restriction of $V''$ a $P$ is 0, the representations irreductibles of $P$ which figurent in the constituants of $V''$ are parmi the same for $V$, and the representations irreductibles of $P$ which figurent in the constituants of $V'$ are the constituants of the restriction of $V$ a $P$. one has therefore $\alpha(V'), \alpha(V'') \geq \alpha(V)$ and $\beta(V'), \beta(V'') \leq \beta(V)$.

This decomposition, and \ref{thm:4.13}, permettent of ameliorer the hypothese of \ref{thm:4.6} en $\beta(W) < \alpha(V)$ and $\beta(W) < \alpha(V')/2$.
\end{remark}

\begin{theorem}\label{thm:4.15}
Let $V$ a representation virtuelle of $W$, without constituant modere. Notons simplement $\varepsilon$ the function $\eta \mapsto \varepsilon(\eta \otimes V, \psi, dx)$ of the group $X$ of the quasi-characters of $K^*$ trivial on $U((1-(1/p))\alpha(V))$.
\begin{enumerate}[label=(\roman*)]
\item The function $\varepsilon$ s'ecrit in such a way unique as a product $\varepsilon = \prod_{i=0}^{p-1} \varepsilon_i$, or $\varepsilon_0$ is constant, $\varepsilon_1$ of the form $\eta \mapsto \eta(\gamma)$, for a certain element $\gamma$ of $K^*$ bien defined $\mod^*\, U((1-(1/p))\alpha(V))$, and or the $\varepsilon_i$, for $i \geq 2$, are homogenes of degree $i$, with values in $\bQ_p/\bZ_p(1)$.
\item Si $2 \leq n \leq p$ and $n(\alpha(V) - \alpha(\eta)) > \alpha(V)$ then $\varepsilon_n(\eta)$ vaut 1.
\item the element $\gamma$ of $K^*$, bien defined $\mod^*\, U((1-1/p)\alpha(V))$, such that $\varepsilon_1(\eta) = \eta(\gamma)$ for $\eta \in X$ precise celui of \eqref{eq:4.6.1}, i.e.~lui is congru $\mod^*\, U(\alpha(V)/2)$.
\end{enumerate}
\end{theorem}

Admettons the existence of a decomposition $\varepsilon = \prod_i \varepsilon_i$ as en (i) and prouvons son unicite and the assertions (ii) and (iii). Let $\gamma_0$ a element of $K^*$ verifiant \eqref{eq:4.2.1}. one has necessarily $\varepsilon_0(\eta) = \varepsilon(V, \psi, dx)$, of or
\[
\varepsilon((\eta-1) \otimes V, \psi) = \eta(\gamma) \cdot \prod_{i=2}^{p-1} \varepsilon_i(\eta) \equiv \eta(\gamma) \pmod{*} \bQ_p/\bZ_p(1).
\]
of after \ref{thm:4.2}, the element $\gamma_1 := \gamma \gamma_0^{-1}$ of $K^*$ is in $U(1)$. one has
\[
\varepsilon((\eta-1) \otimes V, \psi) \cdot \eta(\gamma_0)^{-1} = \eta(\gamma_1) \cdot \prod_{i=2}^{p-1} \varepsilon_i(\eta).
\]
Fixons a integer $n$ between 2 and $p$: $2 \leq n \leq p$. of after \ref{thm:4.9}, the restriction of the function $\varepsilon(\eta \otimes V, \psi, dx): X \to \bC^*$ to the group $X_n$ of the quasi-characters of $K^*$ trivial on $U((1-1/n)\alpha(V))$ is of degree $< n$ (utiliser the form \eqref{eq:1.8.2} of the definition \ref{def:1.8}). The same assertion vaut for the function $\varepsilon((\eta-1) \otimes V, \psi) \cdot \eta(\gamma_0)^{-1}$ which s'en deduces by multiplication by a function of $X$ in $\bC^*$ of degree $< 1$. The function $\varepsilon((\eta-1) \otimes V, \psi) \cdot \eta(\gamma_0)^{-1}$ is en outre with values in $\bQ_p/\bZ_p(1)$, this which permet of lui appliquer \ref{lem:1.9}. Faisant $n=p$, one trouve that it determine uniquement the function $\eta \mapsto \eta(\gamma_1)$ and the functions $\varepsilon_i$. A son tour, the function $\eta \mapsto \eta(\gamma_1)$ determine the elements $\gamma_1$ and $\gamma = \gamma_0 \gamma_1$ of $K^*$ $\mod^*\, U((1-1/n)\alpha(V))$. For $n$ quelconque, one obtains (ii). Enfin, (iii) resulte of the cas $n=2$ of (ii), compare a \ref{thm:4.6}.

Prouvons the assertions of existence.

For $V$ of dimension 1, \ref{prop:1.10} fournit a decomposition \ref{thm:4.15}(i):

\begin{lemma}\label{lem:4.16}
Let $V$ of dimension 1, defined by a character sauvagement ramifie $\chi$ of $K^*$, and $a$ a element of $K^*$ such that $\chi(Ey) = \psi(ay)$ for $v(y) > \alpha(\chi)/p$. the assertion \ref{thm:4.15}(i) is vraie for $V$, and one can prendre $\gamma = a^{-1}$.
\end{lemma}

\begin{lemma}\label{lem:4.17}
Let $L/K$ a extension separable finite of $K$. One the munit of the character additif $\psi \circ \Tr_{L/K}$, and of a quelconque mesure of Haar $dx$. Let $V'$ a representation of $W(K/L)$ and $V$ son induite a $W$. One assumed that the one has $\alpha(V') > 0$ and one sets $\alpha = \Psi_{L/K}(\alpha(V'))$. Supposons \ref{thm:4.15}(i) vrai for $V'$, and posons $\varepsilon_1(\eta) = \eta(\gamma')$, with $\gamma' \in L^*$ and $\gamma = N_{L/K}(\gamma')$. Then the assertions of the theorem \ref{thm:4.15}(i) are vraies for $V$ (and $\psi$), si the one y remplace $\alpha(V)$ by $\alpha$.
\end{lemma}

\begin{proof}
For every quasi-character $\eta$ of $K^*$, one has
\[
\varepsilon(\eta \otimes V, \psi) = \varepsilon(\eta \circ N_{L/K} \otimes V', \psi \circ \Tr_{L/K}).
\]
Si one has $\alpha(\eta) < (1-(1/p))\alpha$, the convexite of $\Psi_{L/K}$ fournit
\[
\alpha(\eta \circ N_{L/K}) \leq \Psi_{L/K}(\alpha(\eta)) \leq (1-(1/p)) \Psi_{L/K}(\alpha) = (1-(1/p)) \alpha(V).
\]
the assertion of the lemma is therefore clear.
\end{proof}

\subsection{}
\label{sec:4.18}

Supposons maintenant that $V$ is a representation quelconque of $W$, without constituant modere. By the variante of the theorem of Brauer datum to the chapitre 2, one ecrit $V = \sum n_i \Ind_{W(K/L_i)}^W(\chi_i)$, or $L_i$ is a finite extension of $K$ and $\chi_i$ a quasi-character of $L_i^*$ verifiant $\alpha(\chi_i) = \Psi_{L_i/K}(\alpha(V))$. Appliquant then \ref{lem:4.16} and \ref{lem:4.17}, one sees that, si the one choisit for each $i$ a element $\gamma_i$ of $L_i^*$ such that one ait $\chi_i(Ey) = \psi \circ \Tr_{L_i/K}(\gamma_i^{-1} y)$ for $y \in L_i$ verifiant $v(y) > \alpha(\chi_i)/p$, then one can prendre $\gamma = \prod N_{L_i/K}(\gamma_i)^{n_i}$ in the theorem \ref{thm:4.15}, relativement a $V$ and $\psi$. C.Q.F.D.

\subsection{}
\label{sec:4.19}

The theorem \ref{thm:4.15} and the section 4.1, permettent therefore of attacher a representation $V$ of $W$ (and to the character additif $\psi$ fixe), a element $\gamma = \gamma(V, \psi)$ of $K^*$, bien defined $\mod^*\, U((1-(1/p))\alpha(V))$, and caracterise by the existence of a decomposition
\begin{equation}\label{eq:4.19.1}
\varepsilon(\eta \otimes V, \psi, dx) = \varepsilon(V, \psi, dx) \cdot \eta(\gamma(V, \psi)) \cdot \prod_{i=2}^{p-1} \varepsilon_i(\eta)
\end{equation}
for each character $\eta$ of $K^*$ trivial on $U((1-(1/p))\alpha(V))$, the functions $\varepsilon_i$ being of the functions homogenes of degree $i$ of the group $X$ of these characters, with values in $\bQ_p/\bZ_p(1)$. This element $\gamma(V, \psi)$ satisfies moreover \eqref{eq:4.1.2}, \eqref{eq:4.2.1} and \eqref{eq:4.6.1} and it is of valuation $a(V) - (\dim V) \cdot n(\psi)$.

That se passe-t-it si the one change character additif $\psi$?

The other characters additifs of $K$ are the characters $\psi_a$, $a$ parcourant $K$, defined by $\psi_a(x) = \psi(ax)$; they are not trivial si $a$ is not zero. Si $X$ is a representation virtuelle of dimension 0 of $W$, one has
\[
\varepsilon(X, \psi_a) = \det X(a) \cdot \varepsilon(X, \psi), \quad \text{(cf.~[1], 5.4)}.
\]
It is easy of en deduire the dependance of $\gamma(V, \psi)$ en $\psi$: one has $\det((\eta-1)V) = \eta^{-\dim V}$ and
\begin{equation}\label{eq:4.19.2}
\gamma(V, \psi_a) = a^{-\dim V} \cdot \gamma(V, \psi).
\end{equation}

\subsection{}
\label{sec:4.20}

Let $\cX$ the dual of Pontrjagin of the group additif of $K$. C'is a space vector of dimension 1 on $K$, for the product $a \cdot \psi = \psi_a$. For each integer $m$, one munit sa power tensor $m$\textsuperscript{eme} of the valuation $v$ for which one has $v(a \psi^{\otimes m}) = v_K(a) + m \cdot n(\psi)$, this for every element $a$ of $K$. The formule \eqref{eq:4.19.1} signifie that the element $\gamma(V) = \gamma(V, \psi) \cdot \psi^{\otimes (-\dim V)}$ of $K \otimes \cX^{\otimes (-\dim V)}$ is independant of the character $\psi$ (not trivial) choisi. one has $v(\gamma(V)) = a(V)$ and
\begin{equation}\label{eq:4.20}
\gamma(V^*) = \gamma(V)^{-1} \cdot \omega^{\otimes \dim V},
\end{equation}
or $\omega$ is the character modere donnant the action of $W$ on $\cX$.

Let $L$ a finite extension of $K$. The morphism $y \mapsto \psi \circ \Tr_{L/K}$ of $\cX$ in $\cX_L$ is $K$-lineaire, and se prolonge en a isomorphism $\iota_{L/K}: K \otimes_K \cX_L \simeq \cX$. One defines the norme $N_{L/K}: \cX_L \to \cX^{[L:K]}$ by the equality
\[
N_{L/K}(b \otimes \psi_L) = N_{L/K}(b) \cdot \psi^{[L:K]} \quad \text{pour } \psi_L \in \cX_L, \; b \in L.
\]
With these notations, one has the formulaire suivant.

\subsection{Formulaire}\label{sec:4.21}

\textbf{(4.21.1) Multiplicativite.} Let $V$, $V'$, $V''$ of the representations of $W$. Supposons $V$ extension of $V'$ by $V''$. one has then
\[
\alpha(V) = \inf(\alpha(V'), \alpha(V''))
\]
and
\[
\gamma(V) = \gamma(V') \cdot \gamma(V'') \pmod{*} U\bigl((1-(1/p))\alpha(V)\bigr).
\]

\textbf{(4.21.2) Induction.} Let $L$ a finite extension of $K$. Let $V'$ a representation of $W(K/L)$ and $V$ son induite a $W$. one has
\[
\alpha(V) = \Psi_{L/K}(\alpha(V'))
\]
and
\[
\gamma(V) = N_{L/K}(\gamma(V')) \pmod{*} U_K\bigl((1-(1/p))\Psi_{L/K}(\alpha(V'))\bigr).
\]

\textbf{(4.21.3) Restriction.} Let $L$ a finite extension of $K$. Let $V$ a representation of $W$ and $V'$ sa restriction a $W(K/L)$. Assume that the one ait $\alpha(L/K) < \alpha(V)$ and notons $e$ the indice of ramification of $L/K$. one has
\[
\gamma(V) = \iota_{L/K}(\gamma(V')) \pmod{*} U_L\bigl((1-(1/p))e(\alpha(V) - \alpha(L/K))\bigr).
\]

\begin{remark}
One notera that si the extension $L/K$ is moderement ramifiee, this formule determine $\gamma(V)$, $\mod^*\, U((1-(1/p))\alpha(V))$, a partir of $\gamma(V')$, and that $\gamma(V)$ ne depend therefore that of the restriction of $V$ to the group of inertie sauvage $P$. This fait can also se deduire of \ref{thm:4.13}(i).
\end{remark}

\textbf{(4.21.4) Dimension 1.} Let $V$ a representation of dimension 1, corresponding to the quasi-character $\chi$ of $K^*$. Supposons $\chi$ sauvagement ramifie, and choisissons a character additif prolongeant the character $\chi \circ \psi(Ey)$ of $A(\alpha(\chi)/p)$. one has
\[
\gamma = \psi \otimes a^{-1} \pmod{*} U((1-(1/p))\alpha(\chi)).
\]

\subsection{}
\label{sec:4.22}

the assertion (4.21.1) resulte aussitot of the definitions, and of the multiplicativite of the factors $\varepsilon_i$. the assertion (4.21.2) reformule \ref{lem:4.17} and (4.21.4) reformule \ref{lem:4.16}. Prouvons (4.21.3). Procedant as en 4.4, 4.8, one can supposer that $V$ is induite of a quasi-character sauvagement ramifie $\chi$ of a finite extension $M$ of $K$, verifiant $\alpha(M/K) \leq \Psi_{M/K}(\alpha(\chi))$ (one pourrait same supposer $\alpha(M/K) < \Psi_{M/K}(\alpha(\chi))$, but it n'importe). Let $t$ a element of $T$ such that $v_M(t) = \alpha(\chi)$. Decomposons $M \otimes_K L$ en a product of field:
\[
M \otimes_K L = \prod_{i \in I} M_i,
\]
and posons $\chi_i = \chi \circ N_{M_i/M}$. one has
\[
V = \Ind_{W(K/M)}^W(\chi), \quad \alpha(V) = v_M(t) - \alpha(\chi) = \Psi_{M/K}(t)
\]
and
\[
V' = \bigoplus_i \Ind_{W(K/M_i)}^{W(K/L)}(\chi_i), \quad \alpha(V') = \sup_i \alpha(\chi_i) = \Psi_{M_i/L}(t).
\]
Indeed, the hypothese $\alpha(L/K) < \alpha(V)$ assure that, for every embedding $\sigma$ of $L$ in $\overline{K}$, $W(K/\sigma(L))$ contient $W[t]$, tandis that the hypothese $\alpha(M/K) \leq \Psi_{M/K}(\alpha(\chi))$ assure that, for every embedding $\tau$ of $M$ in $\overline{K}$ and all $t' > t$, $W(K/\tau(M))$ contient $W[t']$, that $\chi$ is trivial on $W[t']$ and not trivial on $W(K/\tau(M))[t] = W(K/L) \cap W[t]$.

Fixons a character additif not trivial $\psi$ of $K$, and fixons egalement $a$ in $M$ such that one ait $\chi(Ey) = \psi \circ \Tr_{M/K}(ay)$ for $y \in M$ verifiant $v(y) > \alpha(\chi)/p$, and for each $i$ a element $a_i$ of $M_i$ such that one ait $\chi(Ey) = \psi \circ \Tr_{M_i/K}(a_i y)$ for $y \in M_i$ verifiant $v(y) > \alpha(\chi_i)/p$. It s'agit of prouver the congruence
\[
\gamma(V) \equiv \iota_{L/K}(\gamma(V')) \pmod{*} U((1-(1/p))e(\alpha(V) - \alpha(L/K))).
\]
one has by definition $\gamma(V) = N_{M/K}(a^{-1}) \cdot \psi^{-\dim V}$ and $\gamma(V') = \prod_i N_{M_i/L}(a_i^{-1}) \cdot (\psi \circ \Tr_{L/K})^{-\dim V'}$. Or $\dim V = [M:K]$ and $\dim V' = [M:K]/[L:K] \cdot [L:K] = [M:K]$. Moreover $N_{M/K}(a^{-1}) = \prod_i N_{M_i/L}(N_{M_i/M}(a^{-1}))$ and, by the proposition \ref{prop:3.6}, one has
\[
a_i \equiv a \pmod{*} U((1-(1/p))e(\alpha(V) - \alpha(L/K))),
\]
of or the formule (4.21.3).


\begin{thebibliography}{9}
\addcontentsline{toc}{section}{Bibliographie}

\bibitem[1]{D1}
P.~Deligne, \textit{The constants of the functional equations of the functions $L$}, in: Modular Functions of One Variable II, Lecture Notes in Math. 349, Springer-Verlag, 1973, pp.~501--597.

\bibitem[2]{S}
J.-P.~Serre, \textit{Field local}, Hermann, Paris, 1962.

\bibitem[3]{D2}
P.~Deligne, \textit{The local constants of the functional equation of the function $L$ of Artin of a representation orthogonale}, Invent. Math. \textbf{35} (1976), 299--316.

\bibitem[4]{L}
G.~Laumon, \textit{The constants of the functional equations}, Seminar one Arithmetic Bundles: The Mordell Conjecture (Paris, 1980/81), Asterisque \textbf{127} (1985), 7--24.

\end{thebibliography}

\end{document}
